\documentclass{beamer}

% Estilo clásico de Beamer, siguiendo el ejemplo de Overleaf:
% https://es.overleaf.com/learn/latex/Beamer
\usetheme{Madrid}
\usecolortheme{default}
\usefonttheme{professionalfonts}

% Limpieza visual habitual en Beamer
\setbeamertemplate{navigation symbols}{}

\usepackage[utf8]{inputenc}
\usepackage[T1]{fontenc}
\usepackage[spanish,es-noquoting]{babel}
\usepackage{amsmath,amssymb,mathtools,bm}
\usepackage{physics}
\usepackage{tikz}
\usepackage{bm}


\title[Lanczos-Lovelock]{Lagrangianos desde curvatura y métrica}
\subtitle{Sección 2.1: desarrollo paso a paso}
\author[Guillermo Jordán A.]{Guillermo Jordán A.}
\institute[PUCP]{Pontificia Universidad Católica del Perú}
\date{\today}

\newcommand{\Lie}{\mathcal{L}}
\newcommand{\Rcal}{\mathcal{R}}
\renewcommand{\dd}{\mathrm{d}}
\newcommand{\Eab}{E_{ab}}
\newcommand{\Eq}[1]{\eqref{#1}}

% ================================================================
% ROADMAP DE LA DEDUCCIÓN
% ================================================================

% Etapas globales:
% 1 = Identidad principal
% 2 = Variación de la acción
% 3 = Diagnóstico de derivadas superiores
% 4 = Condición de segundo orden
% 5 = Deducción completada
% 0 = Ocultar roadmap

\newcount\roadmapstage
\roadmapstage=0

\newcommand{\SetRoadmapStage}[1]{%
    \global\roadmapstage=#1\relax
}

% ================================================================
% ROADMAP GLOBAL: versión compacta
% ================================================================

\setbeamercolor{roadmap done}{
    use=structure,
    fg=structure.fg,
    bg=structure.fg!12!white
}

\setbeamercolor{roadmap current}{
    use=structure,
    fg=white,
    bg=structure.fg
}

\setbeamercolor{roadmap pending}{
    fg=black!55,
    bg=black!5
}

\setbeamercolor{roadmap page}{
    fg=black!55,
    bg=white
}

% ================================================================
% INDICADOR DEL SUBPASO ACTUAL
% ================================================================

\setbeamercolor{roadmap cue}{
    use=structure,
    fg=structure.fg,
    bg=structure.fg!7!white
}

\newcommand{\RoadmapCue}[1]{%
    \begin{beamercolorbox}[
        wd=\textwidth,
        sep=.55ex,
        leftskip=.8ex,
        rounded=true
    ]{roadmap cue}

    \scriptsize
    \enspace #1

    \end{beamercolorbox}

    \vspace{.5em}
}

%------------------------------------------------
% Caja de cada etapa

\newcommand{\RoadmapBox}[2]{%

    % Etapa terminada
    \ifnum#1<\roadmapstage

        \begin{beamercolorbox}[
            wd=.235\paperwidth,
            ht=2.5ex,
            dp=.8ex,
            center
        ]{roadmap done}

        \tiny\bfseries
        $\checkmark$\; #1\enspace #2

        \end{beamercolorbox}

    % Etapa actual
    \else\ifnum#1=\roadmapstage

        \begin{beamercolorbox}[
            wd=.235\paperwidth,
            ht=2.5ex,
            dp=.8ex,
            center
        ]{roadmap current}

        \tiny\bfseries
        #1\enspace #2

        \end{beamercolorbox}

    % Etapa pendiente
    \else

        \begin{beamercolorbox}[
            wd=.235\paperwidth,
            ht=2.5ex,
            dp=.8ex,
            center
        ]{roadmap pending}

        \tiny\bfseries
        #1\enspace #2

        \end{beamercolorbox}

    \fi\fi
}

%------------------------------------------------
% Footer completo

\setbeamertemplate{footline}{%

    \ifnum\roadmapstage>0

        \leavevmode

        \hbox{%

            \RoadmapBox{1}{Derivadas de Lie}%

            \RoadmapBox{2}{Variación de la acción}%

            \RoadmapBox{3}{Condición de 2° orden}%

            \RoadmapBox{4}{Generalización}%

            \RoadmapBox{5}{Casos}%

            \RoadmapBox{6}{Discusión}%

            \begin{beamercolorbox}[
                wd=.06\paperwidth,
                ht=2.5ex,
                dp=.8ex,
                center
            ]{roadmap page}

            \tiny
            \insertframenumber

            \end{beamercolorbox}%
        }

    \fi
}




% Tabla de contenidos automática al inicio de cada sección
\AtBeginSection[]
{
  \begin{frame}
    \frametitle{Contenido}
    \tableofcontents[currentsection]
  \end{frame}
}

\begin{document}

%================================================
\frame{\titlepage}

%================================================
\begin{frame}
\frametitle{Contenido}
\tableofcontents
\end{frame}

%================================================
\SetRoadmapStage{1}

\begin{frame}
\frametitle{Introducción}
\small

Queremos estudiar una acción gravitacional general (sin considerar la materia):

\[
A=\int_V \dd^Dx\,\sqrt{-g}\,L,
\]

donde el lagrangiano es un escalar construido con la métrica y el tensor de Riemann:

\[
L=L(g_{ab},R_{abcd}).
\]

\begin{block}{Pregunta central}
¿Qué condición debe cumplir $L$ para que las ecuaciones de movimiento no tengan derivadas de la métrica mayores que segundo orden?
\end{block}

La respuesta llevará a exigir:

\[
\nabla_aP^{abcd}=0,
\]

donde $P^{abcd}\equiv \pdv{L}{R_{abcd}}$.

\end{frame}

%================================================
\begin{frame}
\frametitle{Estrategia general}
\small

Seguiremos el orden de la deducción:

\begin{enumerate}
    \item Usar que $L$ es un escalar.
    \item Calcular $\Lie_\xi L$ de dos maneras.
    \item Comparar ambas expresiones.
    \item Obtener la identidad
    \[
    P^{ab}=-2\Rcal^{ab}.
    \]
    \item Variar la acción.
    \item Aislar el tensor de campo $E_{ab}$.
    \item Ver de dónde aparece el término peligroso
    \[
    -2\nabla^m\nabla^nP_{amnb}.
    \]
    \item Imponer la condición de segundo orden.
\end{enumerate}

\end{frame}

%================================================
\section{Derivada de Lie e identidad principal (LL)}

%================================================
\begin{frame}
\frametitle{Derivada de Lie de un tensor general}

\RoadmapCue{
Preparar la herramienta que usaremos para calcular
$\Lie_\xi L$
}
\small

Para un tensor general

\[
T^{a_1\cdots a_r}{}_{b_1\cdots b_s},
\]

la derivada de Lie respecto a un campo vectorial $X^m$ es

\[
(\Lie_XT)^{a_1\cdots a_r}{}_{b_1\cdots b_s}
=
X^m\nabla_mT^{a_1\cdots a_r}{}_{b_1\cdots b_s}
\]

\[
-\sum_{i=1}^{r}
T^{a_1\cdots m\cdots a_r}{}_{b_1\cdots b_s}
\nabla_mX^{a_i}
+
\sum_{j=1}^{s}
T^{a_1\cdots a_r}{}_{b_1\cdots m\cdots b_s}
\nabla_{b_j}X^m.
\]

\begin{block}{Idea}
La derivada de Lie mide cómo cambia un tensor al arrastrarlo a lo largo del flujo generado por $X^m$.
\end{block}

\end{frame}

%================================================
\begin{frame}
\frametitle{Primer cálculo de $\Lie_\xi L$: usando que $L$ es escalar}
\RoadmapCue{
Calcular $\Lie_\xi L$ usando únicamente que $L$ es un escalar
}
\small

Como $L=L(g_{ab},R_{abcd})$ es un escalar,

\[\Lie_\xi L=\xi^m\nabla_m L\]

Entonces, por regla de la cadena covariante,

\[
\nabla_mL
=
\pdv{L}{g_{ab}}\nabla_mg_{ab}
+
\pdv{L}{R_{abcd}}\nabla_mR_{abcd}.
\]

Definimos

\begin{equation}
P^{ab}\equiv \pdv{L}{g_{ab}},
\qquad
P^{abcd}\equiv \pdv{L}{R_{abcd}}.
\label{eq:def-P}
\end{equation}

que se pueden interpretar como la sensibilidad del Lagrangiano ante cambios en la métrica o en la curvatura.

Entonces, \[\nabla_mL=P^{ab}\nabla_mg_{ab}+P^{abcd}\nabla_mR_{abcd}\]

\end{frame}

%================================================
\begin{frame}
\frametitle{Uso de la conexión de Levi-Civita}
\small

En relatividad general usamos la conexión de Levi-Civita, que es compatible con la métrica: $\nabla_mg_{ab}=0$. Por eso, $\nabla_mL=P^{abcd}\nabla_mR_{abcd}$.\\

Así,

\begin{equation}
\boxed{
\Lie_\xi L
=
\xi^mP^{abcd}\nabla_mR_{abcd}.
}
\label{eq:lie-L-primer-calculo}
\end{equation}

\begin{block}{Interpretación}
Este es el primer cálculo de $\Lie_\xi L$. Aquí usamos solamente que $L$ es un escalar y que la conexión preserva la métrica.
\end{block}

Ahora variamos $L$ y, por regla de la cadena, una variación general se escribe como

\[
\delta L
=
\frac{\partial L}{\partial g_{ab}}\delta g_{ab}
+
\frac{\partial L}{\partial R_{ijkl}}\delta R_{ijkl}.
\]

\end{frame}

%================================================
\begin{frame}
\frametitle{Segundo cálculo de $\mathcal L_\xi L$: variando sus argumentos}
\RoadmapCue{
Calcular la misma cantidad variando los argumentos
$g_{ab}$ y $R_{abcd}$
}
\small

Reemplazamos esta última expresión con \Eq{eq:def-P} y entonces

\[
\boxed{
\delta L=P^{ab}\delta g_{ab}+P^{ijkl}\delta R_{ijkl}
}
\]

Luego, un difeomorfismo infinitesimal mueve los puntos de la variedad:

\[
x^a \longrightarrow x^a+\epsilon \xi^a.
\]
Sea una función escalar y arbitraria $f$. La expandimos por Taylor:

\[
f(x+\epsilon\xi) = f(x) + \epsilon \xi^a \partial_a f + O(\epsilon^2)
\]
Su variación será:
\[
\delta f(x) = f(x+\epsilon\xi) - f(x) = \epsilon \xi^a \partial_a f + O(\epsilon^2) \approx \epsilon \xi^a \partial_a f
\]
Pero sabemos que $\mathcal{L}_\xi f = \xi^a \partial_a f$. Por lo tanto, para escalares, 
\[ \delta f(x) = \mathcal{L}_\xi f \]

\end{frame}

%================================================
\begin{frame}
\frametitle{Por qué aparece la derivada de Lie}
\small

Por eso, para la variación inducida por el difeomorfismo,

\[
\delta_\xi g_{ab}=\mathcal L_\xi g_{ab},
\qquad
\delta_\xi R_{ijkl}=\mathcal L_\xi R_{ijkl}.
\]

Así,

\begin{equation}
\boxed{
\mathcal L_\xi L
=
P^{ab}\mathcal L_\xi g_{ab}
+
P^{ijkl}\mathcal L_\xi R_{ijkl}.
}
\label{eq:lie-L-regla-cadena}
\end{equation}

Para la métrica,

\[
\mathcal L_\xi g_{ab}
=
\xi^m\nabla_mg_{ab}
+
g_{mb}\nabla_a\xi^m
+
g_{am}\nabla_b\xi^m.
\]

Como usamos la conexión de Levi-Civita, $\nabla_mg_{ab}=0$. Entonces,

\[
\mathcal L_\xi g_{ab}
=
g_{mb}\nabla_a\xi^m
+
g_{am}\nabla_b\xi^m.
\]

Bajando el índice llegamos a,

\begin{equation}
\boxed{
\mathcal L_\xi g_{ab}
=
\nabla_a\xi_b+\nabla_b\xi_a.
}
\label{eq:lie-metrica}
\end{equation}

\end{frame}

%================================================
\begin{frame}
\frametitle{Derivada de Lie del tensor de Riemann}
\small

Para la curvatura $R_{ijkl}$, por definición, la derivada de Lie de un tensor completamente covariante $\mathcal L_\xi R_{ijkl}$ es

\begin{equation}\boxed{
\xi^m\nabla_mR_{ijkl}
+
R_{mjkl}\nabla_i\xi^m
+
R_{imkl}\nabla_j\xi^m
+
R_{ijml}\nabla_k\xi^m
+
R_{ijkm}\nabla_l\xi^m
}
\label{eq:lie-riemann}
\end{equation}

Sustituyendo \Eq{eq:lie-metrica} y \Eq{eq:lie-riemann} en \Eq{eq:lie-L-regla-cadena}, llegamos a:

\[
\mathcal L_\xi L
=
P^{ab}(\nabla_a\xi_b+\nabla_b\xi_a)
\]\[+P^{ijkl}
\left(
\xi^m\nabla_mR_{ijkl}
+
R_{mjkl}\nabla_i\xi^m
+
R_{imkl}\nabla_j\xi^m
+
R_{ijml}\nabla_k\xi^m
+
R_{ijkm}\nabla_l\xi^m
\right)
\]

A partir de Ahora simplificamos por partes: primero el \textbf{término métrico}, luego los cuatro \textbf{términos de curvatura}. \newline

La parte métrica es $P^{ab}(\nabla_a\xi_b+\nabla_b\xi_a)$ que separamos como $P^{ab}\nabla_a\xi_b+P^{ab}\nabla_b\xi_a$.

\end{frame}

%================================================
\begin{frame}
\frametitle{Simplificación del término métrico}
\small

En el segundo término intercambiamos índices mudos: $a\leftrightarrow b$. Entonces $P^{ab}\nabla_b\xi_a
=P^{ba}\nabla_a\xi_b$. \newline

Como $P^{ab}=P^{ba}$, por su simetría, tenemos $P^{ba}\nabla_a\xi_b = P^{ab}\nabla_a\xi_b$. \newline

Por tanto,

\begin{equation}
\boxed{P^{ab}\mathcal L_\xi g_{ab} = 2P^{ab}\nabla_a\xi_b}
\label{eq:contraccion-P-lie-metrica}
\end{equation}

La parte con el tensor de Riemann es $P^{ijkl}\mathcal L_\xi R_{ijkl}$. Expandiendo,

\[
P^{ijkl}\mathcal L_\xi R_{ijkl}
=
P^{ijkl}\xi^m\nabla_mR_{ijkl}
\]\[
+
P^{ijkl}R_{mjkl}\nabla_i\xi^m
+
P^{ijkl}R_{imkl}\nabla_j\xi^m
+
P^{ijkl}R_{ijml}\nabla_k\xi^m
+
P^{ijkl}R_{ijkm}\nabla_l\xi^m.
\]

\begin{block}{Nota}
El primer término ya coincide con el que apareció cuando tratamos a $L$ como escalar. Los otros cuatro términos todavía deben reorganizarse. Eso se hará a partir de las simetrías algebraicas entre $R_{ijkl}$ y $P^{ijkl}$.
\end{block}

\end{frame}

\begin{frame}{¿Por qué \(P^{ijkl}\) hereda las simetrías de \(R_{ijkl}\)?}
\RoadmapCue{
Justificar las simetrías necesarias para simplificar
el término de curvatura
}
\small
Definimos \(P^{ijkl}\) mediante $\delta L=P^{ijkl}\delta R_{ijkl}$. Pero \(R_{ijkl}\) no varía libremente. Como
\[
R_{ijkl}=-R_{jikl},\qquad
R_{ijkl}=-R_{ijlk},\qquad
R_{ijkl}=R_{klij},
\]
entonces también
\[
\delta R_{ijkl}=-\delta R_{jikl},\qquad
\delta R_{ijkl}=-\delta R_{ijlk},\qquad
\delta R_{ijkl}=\delta R_{klij}.
\]

Ahora descomponemos \(P^{ijkl}\) en el primer par:
\[
P^{ijkl}=P^{(ij)kl}+P^{[ij]kl}.
\]

La parte simétrica no contribuye: $X=P^{(ij)kl}\delta R_{ijkl}$. Intercambiando los índices mudos \(i\leftrightarrow j\),
\[
X=P^{(ji)kl}\delta R_{jikl}
  =P^{(ij)kl}(-\delta R_{ijkl})
  =-X.
\]
Por tanto, $X=0$. Así, $P^{ijkl}\delta R_{ijkl} = P^{[ij]kl}\delta R_{ijkl}$.

\[
\Rightarrow \boxed{P^{ijkl}=-P^{jikl}}.
\]

\end{frame}

\begin{frame}{¿Por qué \(P^{ijkl}\) hereda las simetrías de \(R_{ijkl}\)?}

El mismo argumento se aplica al segundo par. Como $\delta R_{ijkl}=-\delta R_{ijlk}$, solo contribuye la parte antisimétrica de \(P^{ijkl}\) en \(k,l\). Entonces
\[
\boxed{P^{ijkl}=-P^{ijlk}}.
\]

Además, como $\delta R_{ijkl}=\delta R_{klij}$, solo contribuye la parte simétrica de \(P^{ijkl}\) bajo el intercambio de pares:
\[
(ij)\leftrightarrow(kl).
\]

Por tanto,
\[
\boxed{P^{ijkl}=P^{klij}}.
\]

Es decir, \(P^{ijkl}\) puede tomarse con las mismas simetrías algebraicas que \(R_{ijkl}\), porque cualquier parte de \(P\) que no tenga esas simetrías no contribuye a
\[
\delta L=P^{ijkl}\delta R_{ijkl}.
\]

\end{frame}

%================================================
\begin{frame}
\frametitle{Por qué los cuatro términos son iguales}
\small

Usamos que $P^{ijkl}$ tiene las mismas simetrías algebraicas que $R_{ijkl}$:

\[
R_{ijkl}=-R_{jikl},
\qquad
R_{ijkl}=-R_{ijlk},
\qquad
R_{ijkl}=R_{klij}
\]\[P^{ijkl}=-P^{jikl},
\qquad
P^{ijkl}=-P^{ijlk},
\qquad
P^{ijkl}=P^{klij}
\]


Con estas simetrías e intercambiando índices mudos, los cuatro términos con $\nabla\xi$ se convierten en cuatro copias de uno solo: $P^{ijkl}R_{mjkl}\nabla_i\xi^m$. \newline

Por tanto, \[\boxed{P^{ijkl}\mathcal L_\xi R_{ijkl} = P^{ijkl}\xi^m\nabla_mR_{ijkl}+4P^{ijkl}R_{mjkl}\nabla_i\xi^m}\]

Antes de continuar, mostraremos el porqué se heredan dichas simetrías algebraicas entre $R_{ijkl}$ y $P^{ijkl}$.

\end{frame}

%================================================
\begin{frame}
\frametitle{Aparición de $\mathcal R^{ab}$}
\small

\begin{block}{Nota}
\(P^{abcd}=\left(\partial L/\partial R_{abcd}\right)_{g_{ij}}\) mide la respuesta del lagrangiano gravitacional frente a la curvatura y actúa como la densidad local de entropía de Wald . Revisar \cite{Padmanabhan2010}.
\end{block}

Definimos el tensor

\begin{equation} 
\boxed{\mathcal R^{ab} \equiv P^{aijk}R^b{}_{ijk}}
\label{eq:def-Rcal}
\end{equation}

Ahora, el término $P^{ijkl}R_{mjkl}\nabla_i\xi^m$ de puede escribirse, renombrando índices y bajando el índice de $\xi^m$, como

\[
\mathcal R^{ab}\nabla_a\xi_b.
\]

Entonces,

\begin{equation}\boxed{
P^{ijkl}\mathcal L_\xi R_{ijkl} = 
P^{ijkl}\xi^m\nabla_mR_{ijkl}+4\mathcal R^{ab}\nabla_a\xi_b}
\label{eq:contraccion-P-lie-riemann}
\end{equation}

\end{frame}

%================================================
\begin{frame}
\frametitle{Resultado del segundo cálculo}
\small

De \Eq{eq:def-Rcal} y \Eq{eq:contraccion-P-lie-riemann} en \Eq{eq:lie-L-regla-cadena} tenemos que

\[
\mathcal L_\xi L
=
2P^{ab}\nabla_a\xi_b
+
P^{ijkl}\xi^m\nabla_mR_{ijkl}
+
4\mathcal R^{ab}\nabla_a\xi_b.
\]

Agrupamos los términos con $\nabla_a\xi_b$:

\[
\mathcal L_\xi L
=
P^{ijkl}\xi^m\nabla_mR_{ijkl}
+
2\left(P^{ab}+2\mathcal R^{ab}\right)\nabla_a\xi_b.
\]

Por tanto,

\begin{equation}
\boxed{
\mathcal L_\xi L
=
P^{ijkl}\xi^m\nabla_mR_{ijkl}
+
2\left(P^{ab}+2\mathcal R^{ab}\right)\nabla_a\xi_b.
}
\label{eq:lie-L-segundo-calculo}
\end{equation}

Ahora recordemos que en \Eq{eq:lie-L-primer-calculo}, obtuvimos

\[
\mathcal L_\xi L
=
\xi^mP^{ijkl}\nabla_mR_{ijkl}
\]
Igualamos \Eq{eq:lie-L-primer-calculo} y \Eq{eq:lie-L-segundo-calculo}.

\end{frame}

%================================================
\begin{frame}
\frametitle{Comparación con el cálculo escalar}
\RoadmapCue{
Comparar las dos expresiones obtenidas para la misma cantidad
$\Lie_\xi L$
}
\scriptsize

Obtenemos que el coeficiente de $\nabla_a\xi_b$ debe anularse:

\[
P^{ab}+2\mathcal R^{ab}=0.
\]

Finalmente,

\begin{equation}
\boxed{
P^{ab}=-2\mathcal R^{ab}
}
\label{eq:identidad-P-Rcal}
\end{equation}

Ahora, notemos que como la métrica es simétrica,
\[
g_{ab}=g_{ba}
\quad \Rightarrow \quad
\delta g_{ab}=\delta g_{ba}.
\]

Descomponemos
\[
P^{ab}=P^{(ab)}+P^{[ab]}.
\]

La parte antisimétrica no contribuye:
\[
X=P^{[ab]}\delta g_{ab}.
\]

Intercambiando índices mudos \(a\leftrightarrow b\),
\[
X=P^{[ba]}\delta g_{ba}
=
(-P^{[ab]})\delta g_{ab}
=
-X.
\]

\end{frame}


%================================================
\begin{frame}
\frametitle{Consecuencia importante}

Por tanto, $X=0 \quad \Rightarrow \quad P^{ab}\delta g_{ab}=P^{(ab)}\delta g_{ab}$.

Así, la derivada respecto a una métrica simétrica puede tomarse simétrica:
\[
\boxed{P^{ab}=P^{ba}}
\]

\begin{block}{Simetria en $\Rcal^{ab}$}
Usando la identidad en \Eq{eq:identidad-P-Rcal} se obtiene inmediatamente que $\Rcal^{ab}$ es simétrico, porque $P^{ab}$ lo es:

\begin{equation}
P^{ab}=P^{ba}
\quad\Longrightarrow\quad
\Rcal^{ab}=\Rcal^{ba}
\label{eq:simetria-Rcal}
\end{equation}

\end{block}

\RoadmapCue{
Ahora aplicaremos la teoría vista hasta acá para el caso más elemental de Hilbert-Einstein.
}

\end{frame}

%================================================
\begin{frame}
\frametitle{Caso Einstein--Hilbert: cálculo de \(P^{abcd}\)}
\scriptsize

Para Einstein--Hilbert tomamos $L=R$.

Recordemos que el escalar de Ricci se obtiene contrayendo el tensor de Riemann:
\begin{equation}
R=g^{ac}g^{bd}R_{abcd} = g^{ad}g^{bc}R_{abdc}
\label{eq:ricci-contraccion-riemann}
\end{equation}

Por definición,
\[
P^{abcd}\equiv
\left(
\frac{\partial L}{\partial R_{abcd}}
\right)_{g_{ij}}.
\]

Entonces, para \(L=R\) uno esperaría \(g^{ac}g^{bd}\), pero aplicamos la expresión en \Eq{eq:ricci-contraccion-riemann}:
\[
P^{abcd}_{EH}
= \frac{\partial L}{\partial R} \times
\frac{\partial R}{\partial R_{abcd}} = \frac{\partial R}{\partial R_{abcd}} = \frac{\partial (g^{ac}g^{bd}R_{abcd})}{\partial R_{abcd}} = \frac{1}{2}\frac{\partial \left( g^{ac}g^{bd}R_{abcd} + g^{ad}g^{bc}R_{abdc} \right)}{\partial R_{abcd}} 
\]

Aplicando la antisimetría del tensor de Riemann:
\[
R_{abcd}=-R_{abdc} \qquad \longrightarrow \qquad P^{abcd}_{EH} = \frac{1}{2}\frac{\partial \left( g^{ac}g^{bd}R_{abcd} - g^{ad}g^{bc}R_{abcd} \right)}{\partial R_{abcd}} 
\]

Que finalmente nos da:
\begin{equation}
P^{abcd}_{EH} =\frac12 \left( g^{ac}g^{bd} - g^{ad}g^{bc} \right)
\label{eq:P-einstein-hilbert}
\end{equation}

Este tensor cumple las simetrías de \(R_{abcd}\):
\[
P^{abcd}_{EH}=-P^{abdc}_{EH},
\qquad
P^{abcd}_{EH}=P^{cdab}_{EH}.
\]

\end{frame}

%================================================
\begin{frame}
\frametitle{Cálculo de \(\mathcal R_{ab}\) para Einstein--Hilbert}

El tensor de Ricci generalizado se define como $\mathcal R_{ab}\equiv P_a{}^{ijk}R_{bijk}$. \newline

Para Einstein--Hilbert,
\[
P^{mijk}_{EH}
=
\frac12
\left(g^{mj}g^{ik}-g^{mk}g^{ij}\right)
\]

Bajamos el primer índice: $P_a{}^{ijk} = g_{am}P^{mijk}$. \newline

Entonces
\[
\begin{aligned}
P_a{}^{ijk}&=g_{am}
\frac12
\left(g^{mj}g^{ik}-g^{mk}g^{ij}\right)
=
\frac12
\left(\delta_a^j g^{ik}-\delta_a^k g^{ij}\right)
\end{aligned}
\]

Ahora reemplazamos en \(\mathcal R_{ab}\):
\[
\begin{aligned}
\mathcal R_{ab}=P_a{}^{ijk}R_{bijk}=
\frac12
\left(\delta_a^j g^{ik}-\delta_a^k g^{ij}\right)
R_{bijk}=
\frac12
\left(g^{ik}R_{biak}-g^{ij}R_{bija}\right)
\end{aligned}
\]

\end{frame}

%================================================
\begin{frame}
\frametitle{Por qué \(\mathcal R_{ab}=R_{ab}\)}
\scriptsize

Tenemos
\[
\mathcal R_{ab}
=\frac12
\left(g^{ik}R_{biak}-g^{ij}R_{bija}\right)
\]

Primer término: $g^{ik}R_{biak}$. \newline

Usando intercambio de pares $R_{biak}=R_{akbi}$. Luego, usando antisimetría en el segundo par, $R_{akbi}=-R_{akib}$. Y usando antisimetría en el primer par, $-R_{akib}=R_{kaib}$.

Por tanto,
\[
g^{ik}R_{biak}
=g^{ik}R_{akbi} = -g^{ik}R_{akib}
=g^{ik}R_{kaib} = g^{ik}R_{iakb}
=R_{ab}
\]

Segundo término: $g^{ij}R_{bija}$ \newline

Usando antisimetría en el primer par, $R_{bija}=-R_{ibja}$. Entonces,
\[
g^{ij}R_{bija}
=
-g^{ij}R_{ibja}
=
-R_{ba}
\]

Como el tensor de Ricci es simétrico, $R_{ba}=R_{ab}$. Así,
\[
g^{ij}R_{bija}=-R_{ab}.
\]

Finalmente,
\[
\mathcal R_{ab}
=\frac12
\left(R_{ab}-(-R_{ab})\right)
=R_{ab}
\]

\[
\boxed{\mathcal R_{ab}=R_{ab}.}
\]

\end{frame}



%================================================
\begin{frame}
\frametitle{Constante cosmológica: momentos y tensor $\mathcal{R}_{ab}$}
\small

Para el término de constante cosmológica,

\[
L_{\Lambda}=-2\Lambda,
\]

no existe dependencia respecto del tensor de Riemann. Por tanto,

\[
P_{\Lambda}^{abcd}
\equiv
\frac{\partial L_{\Lambda}}{\partial R_{abcd}}
=0.
\]

Asimismo, al ser \(L_{\Lambda}\) una constante,

\[
P_{\Lambda}^{ab}
\equiv
\left(
\frac{\partial L_{\Lambda}}{\partial g_{ab}}
\right)_{R_{cdef}}
=0.
\]

El tensor análogo al tensor de Ricci queda entonces

\[
\mathcal{R}_{\Lambda}^{ab}
\equiv
P_{\Lambda}^{aijk}R^{b}{}_{ijk}
=0.
\]

Por consiguiente,

\[
\boxed{
P_{\Lambda}^{abcd}=0,
\qquad
\mathcal{R}_{\Lambda}^{ab}=0
}.
\]

\end{frame}
%================================================



%================================================
\begin{frame}
\frametitle{Resumen de la sección}

\RoadmapCue{
Hemos obtenido las identidades fundamentales que utilizaremos
para variar la acción.
}

\small

A partir de la invariancia bajo difeomorfismos demostramos que:

\[
P^{ab}=-2\mathcal R^{ab} = -2P^{aijk}R^{b}_{ijk}
\]

\vspace{0.2cm}

Además, concluimos que:

\[
P^{ab}=P^{ba}
\qquad\Longleftrightarrow\qquad
\mathcal R^{ab}=\mathcal R^{ba}.
\]

\vspace{0.2cm}

Para el caso de Einstein--Hilbert (\(L=R\)) verificamos que

\[
P^{abcd}
=
\frac12
\left(
g^{ac}g^{bd} - g^{ad}g^{bc}
\right),
\qquad
\mathcal R_{ab}=R_{ab}
\]

Estas identidades muestran que el formalismo reproduce exactamente la Relatividad General en el caso más simple y establecen la base para variar una acción gravitacional completamente general.

\end{frame}


%================================================
\section{Variación de la acción (LL)}
\SetRoadmapStage{2}
%================================================
\begin{frame}
\frametitle{Variación de la acción}
\small

Ahora variamos

\[
A=\int_V\dd^Dx\,\sqrt{-g}L.
\]

Entonces

\[
\delta A
=
\int_V\dd^Dx\,\delta(\sqrt{-g}L).
\]

Usamos

\[
\delta(\sqrt{-g}L)
=
\delta(\sqrt{-g})L+\sqrt{-g}\delta L.
\]

Para variaciones respecto a $g^{ab}$,

\[
\delta\sqrt{-g}
=
-\frac12\sqrt{-g}\,g_{ab}\delta g^{ab}.
\]

Por tanto,

\begin{equation}
\delta(\sqrt{-g}L)
=
\sqrt{-g}
\left[
-\frac12g_{ab}L\delta g^{ab}
+
\delta L
\right]
\label{eq:variacion-raiz-g-L}
\end{equation}

\end{frame}

\begin{frame}{Cambio de variable: \(g_{ab}\) a \(g^{ab}\)}

Antes de seguir, conviene desarrollar las expresiones que ya teníamos, pero con la version contravariante de la métrica.

\begin{block}{Identidad para la métrica contravariante}

Partimos de la identidad ya conocida en \Eq{eq:identidad-P-Rcal}: $\left(\frac{\partial L}{\partial g_{ab}}\right)_{R} = -2\mathcal{R}^{ab}$.

Como $g^{ac}g_{cb}=\delta^a_b$, al variar se obtiene
\begin{equation}
\delta g^{ac}g_{cb} +g^{ac}\delta g_{cb} = 0 \qquad \Rightarrow \qquad
\delta g_{ab}
=
-g_{ac}g_{bd}\,\delta g^{cd}
\label{eq:variacion-metrica-inversa}
\end{equation}

Entonces $\delta_g L=-2\mathcal{R}^{ab}\delta g_{ab} = 2\mathcal{R}^{ab}g_{ac}g_{bd}\delta g^{cd}$.\newline

Pero $\mathcal{R}_{cd}=g_{ac}g_{bd}\mathcal{R}^{ab}$, por lo que $\delta_g L=2\mathcal{R}_{cd}\delta g^{cd}$. Así,
\begin{equation}
\boxed{
\left(\frac{\partial L}{\partial g^{ab}}\right)_{R}
=
2\mathcal{R}_{ab}
}
\label{eq:derivada-L-metrica-contravariante}
\end{equation}
\end{block}

\end{frame}

%================================================
\begin{frame}
\frametitle{Variación del lagrangiano}
\small

Notemos que en $\delta \sqrt{-g}$ aparece $\delta g^{ab}$, por lo que es mucho mas natural tomar $L=L(g^{ab},R_{abcd})$ con la componente contravariante de la metrica. Entonces, variamos

\[
\delta L
=
\underbrace{\left( \dfrac{\partial L}{\partial g^{ab}}
\right)_{R_{ijkl}}
\delta g^{ab}}_{\delta L_g}
+
\underbrace{P^{abcd}\delta R_{abcd}}_{\delta L_R}\]

Usando la identidad \Eq{eq:derivada-L-metrica-contravariante} podemos llegar a que 

\[
\delta L
=
2\Rcal_{ab}\delta g^{ab}
+
P^{abcd}\delta R_{abcd}
\]

Luego reemplazamos esto en \Eq{eq:variacion-raiz-g-L} y llegamos a

\begin{equation}
\delta(\sqrt{-g}L)
=
\sqrt{-g}
\left[
\left(2\Rcal_{ab}-\frac12g_{ab}L\right)\delta g^{ab}
+ P^{abcd}\delta R_{abcd}
\right]
\label{eq:variacion-raiz-g-L-expandida}
\end{equation}

\end{frame}

%================================================
\begin{frame}
\frametitle{El término problemático}
\scriptsize
\RoadmapCue{
\[
\delta(\sqrt{-g}L)
=
\sqrt{-g}
\left[
\left(2\Rcal_{ab}-\frac12g_{ab}L\right)\delta g^{ab}
+ \bm{P^{abcd}\delta R_{abcd}}
\right]
\]
}

El término delicado es $P^{abcd}\delta R_{abcd}$. \newline

Como $R_{abcd}=g_{ae}R^e{}_{bcd}$, su variación es
\[
\delta R_{abcd}
=
\delta g_{ae}R^e{}_{bcd}
+
g_{ae}\delta R^e{}_{bcd}.
\]

Entonces,

\begin{equation}
P^{abcd}\delta R_{abcd}
=
P^{abcd}\delta g_{ae}R^e{}_{bcd}
+
P^{abcd}g_{ae}\delta R^e{}_{bcd}
\label{eq:separacion-variacion-riemann}
\end{equation}

El primer aporte es $P^{abcd}\delta g_{ae}R^e{}_{bcd}$. Reemplazando con la identidad en \Eq{eq:variacion-metrica-inversa} definida como $ \delta g_{ae}=-g_{am}g_{en}\delta g^{mn}$ obtenemos

\[
P^{abcd}R^e{}_{bcd}\delta g_{ae}
=
-P^{abcd}R^e{}_{bcd}g_{am}g_{en}\delta g^{mn} = -\Rcal_{mn}\delta g^{mn}
\]

Por tanto,

\begin{equation}
\boxed{
P^{abcd}\delta R_{abcd} = -\Rcal_{ab}\delta g^{ab} + P^{abcd}g_{ae}\delta R^e{}_{bcd}
}
\label{eq:variacion-riemann-Rcal}
\end{equation}

\end{frame}

%================================================
\begin{frame}
\frametitle{Aplicación de la identidad de Palatini}
\scriptsize
\RoadmapCue{
\[
\delta(\sqrt{-g}L)
=
\sqrt{-g}
\left[
\left(2\Rcal_{ab}-\frac12g_{ab}L\right)\delta g^{ab}
+ \bm{P^{abcd}\delta R_{abcd}}
\right]
\]
}

La variación del tensor de Riemann, por la identidad de Palatini, se escribe como

\[
\boxed{
\delta R^a{}_{bcd}
=
\nabla_c\delta\Gamma^a{}_{db}
-
\nabla_d\delta\Gamma^a{}_{cb}.
}
\]

Además,

\[
\delta\Gamma^a{}_{bc}
=
\frac12 g^{ad}
\left(
\nabla_b\delta g_{cd}
+
\nabla_c\delta g_{bd}
-
\nabla_d\delta g_{bc}
\right).
\]

\begin{block}{Idea}
$\delta R$ contiene derivadas de $\delta \Gamma$, y $\delta\Gamma$ contiene derivadas de $\delta g$. Por eso aparecerán términos de borde al integrar por partes.
\end{block}

Del segundo término de la expresión \Eq{eq:variacion-riemann-Rcal} tenemos $P^{abcd}g_{ae}\delta R^e{}_{bcd}$.\newline

Usando la identidad de Palatini:

\[
P^{abcd}g_{ae}\delta R^e{}_{bcd}
=
P^{abcd}g_{ae}
\left(
\nabla_c\delta\Gamma^e{}_{db}
-
\nabla_d\delta\Gamma^e{}_{cb}
\right)
\]

\end{frame}

%================================================
\begin{frame}
\frametitle{Reemplazo de \(\delta\Gamma\) en \(P^{abcd}\delta R_{abcd}\)}
\scriptsize

Partimos de la identidad de Palatini:
\[
\delta R^{e}{}_{bcd}
=
\nabla_c\delta\Gamma^{e}{}_{db}
-
\nabla_d\delta\Gamma^{e}{}_{cb}.
\]

Entonces
\[
P^{abcd}g_{ae}\delta R^{e}{}_{bcd}
=
P^{abcd}g_{ae}
\left(
\nabla_c\delta\Gamma^{e}{}_{db}
-
\nabla_d\delta\Gamma^{e}{}_{cb}
\right).
\]

Usando que \(P^{abcd}\) es antisimétrico en \(c,d\): $P^{abcd}=-P^{abdc}$, los dos términos son iguales y queda
\[
P^{abcd}g_{ae}\delta R^{e}{}_{bcd}
=
2P^{abcd}g_{ae}\nabla_c\delta\Gamma^{e}{}_{db}.
\]

Ahora usamos
\[
\delta\Gamma^{e}{}_{db}
=
\frac12 g^{ei}
\left(
\nabla_d\delta g_{bi}
+
\nabla_b\delta g_{di}
-
\nabla_i\delta g_{db}
\right).
\]

Como \(\nabla_c g_{ae}=0\) y \(g_{ae}g^{ei}=\delta^i_a\),
\[
\begin{aligned}
2P^{abcd}g_{ae}\nabla_c\delta\Gamma^{e}{}_{db}
&=
P^{ibcd}\nabla_c
\left(
\nabla_d\delta g_{bi}
+
\nabla_b\delta g_{di}
-
\nabla_i\delta g_{db}
\right)
\\
&=
P^{ibcd}
\left(
\nabla_c\nabla_d\delta g_{bi}
+
\nabla_c\nabla_b\delta g_{di}
-
\nabla_c\nabla_i\delta g_{db}
\right).
\end{aligned}
\]

Usando las simetrías de \(P^{abcd}\), la simetría
\(\delta g_{ab}=\delta g_{ba}\), y renombrando índices mudos, esta combinación se reduce a
\begin{equation}
\boxed{
P^{abcd}g_{ae}\delta R^{e}{}_{bcd}
=
2P^{ibjd}\nabla_j\nabla_b\delta g_{di}.
}
\label{eq:doble-derivada-variacion-metrica}
\end{equation}

\end{frame}

%================================================
\begin{frame}
\frametitle{Término con \(\delta g_{di}\)}
\scriptsize
\RoadmapCue{
\[
\delta(\sqrt{-g}L)
=
\sqrt{-g}
\left[
\left(2\Rcal_{ab}-\frac12g_{ab}L\right)\delta g^{ab}
+ \bm{P^{abcd}\delta R_{abcd}}
\right]
\]
}

Reescribimos el término a la derecha respecto a \(\nabla_j\):
\[
\begin{aligned}
2P^{ibjd}\nabla_j\nabla_b\delta g_{di}
&=
\nabla_j\left(
2P^{ibjd}\nabla_b\delta g_{di}
\right)
-
2(\nabla_jP^{ibjd})\nabla_b\delta g_{di}.
\end{aligned}
\]

En el segundo término renombramos índices mudos para escribir
\[
(\nabla_jP^{ibjd})\nabla_b\delta g_{di}
=
(\nabla_cP^{ijcd})\nabla_j\delta g_{di}.
\]

Entonces
\[
-2(\nabla_jP^{ibjd})\nabla_b\delta g_{di}
=
-2(\nabla_cP^{ijcd})\nabla_j\delta g_{di}.
\]

Integramos nuevamente por partes, ahora respecto a \(\nabla_j\):
\[
\begin{aligned}
-2(\nabla_cP^{ijcd})\nabla_j\delta g_{di}
&=
-\nabla_j
\left(
2\delta g_{di}\nabla_cP^{ijcd}
\right)
+
2\delta g_{di}\nabla_j\nabla_cP^{ijcd}.
\end{aligned}
\]

Renombrando índices mudos en el último término,
\[
2\delta g_{di}\nabla_j\nabla_cP^{ijcd}
=
2\delta g_{di}\nabla_b\nabla_jP^{ibjd}.
\]
\end{frame}


\begin{frame}{Término con \(\delta g_{di}\)}
\scriptsize
\RoadmapCue{
\[
\delta(\sqrt{-g}L)
=
\sqrt{-g}
\left[
\left(2\Rcal_{ab}-\frac12g_{ab}L\right)\delta g^{ab}
+ \bm{P^{abcd}\delta R_{abcd}}
\right]
\]
}

Por tanto,
\[
\begin{aligned}
2P^{ibjd}\nabla_j\nabla_b\delta g_{di}
&=
\nabla_j
\left(
2P^{ibjd}\nabla_b\delta g_{di}
\right)
-
\nabla_j
\left(
2\delta g_{di}\nabla_cP^{ijcd}
\right)
\\
&\quad
+
2\delta g_{di}\nabla_b\nabla_jP^{ibjd}
\end{aligned}
\]

Definimos
\begin{equation}
\boxed{
\delta v^j
\equiv
2P^{ibjd}\nabla_b\delta g_{di}
-
2\delta g_{di}\nabla_cP^{ijcd}
}
\label{eq:def-delta-v}
\end{equation}

Así,
\[
2P^{ibjd}\nabla_j\nabla_b\delta g_{di}
=
2\delta g_{di}\nabla_b\nabla_jP^{ibjd}
+
\nabla_j\delta v^j
\]

Finalmente, usando lo anterior y que $\delta g_{di}=-g_{da}g_{ib}\delta g^{ab}$, en \Eq{eq:doble-derivada-variacion-metrica} se obtiene
\begin{equation}
\boxed{
P^{abcd}g_{ae}\delta R^{e}{}_{bcd}
=
-2\nabla^m\nabla^nP_{amnb}\,\delta g^{ab}
+
\nabla_j\delta v^j.
}
\label{eq:termino-curvatura-borde}
\end{equation}

Parra cerrar esta parte, reemplazamos \Eq{eq:termino-curvatura-borde} en \Eq{eq:variacion-riemann-Rcal}:

\begin{equation}
\boxed{
P^{abcd}\delta R_{abcd}
=
-\Rcal_{ab}\delta g^{ab}-2\nabla^m\nabla^nP_{amnb}\,\delta g^{ab} + \nabla_j\delta v^j
}
\label{eq:variacion-riemann-completa}
\end{equation}

\end{frame}

%================================================
\begin{frame}
\frametitle{Término de borde}
\small

El término de borde tiene la forma

\[
\boxed{
\delta v^j
=
2P^{ibjd}\nabla_b\delta g_{di}
-
2\delta g_{di}\nabla_cP^{ijcd}
}
\]

Entonces,

\[
\int_V\dd^Dx\,\sqrt{-g}\nabla_j\delta v^j
\]

se transforma en una integral sobre la frontera.

\begin{alertblock}{Punto importante}
Aunque fijemos
\[
\delta g_{ab}=0
\]
en la frontera, todavía puede quedar
\[
\nabla_b\delta g_{ab}.
\]
Por eso en gravedad se necesitan términos de borde adecuados, igual que en Einstein-Hilbert con Gibbons-Hawking-York.
\end{alertblock}

\end{frame}

%================================================
\begin{frame}
\frametitle{Sustitución en la variación total}
\small

Teníamos

\[
\delta(\sqrt{-g}L)
=
\sqrt{-g}
\left[
\left(2\Rcal_{ab}-\frac12g_{ab}L\right)\delta g^{ab}
+
P^{abcd}\delta R_{abcd}
\right]
\]

Usamos la relación \Eq{eq:variacion-riemann-completa}, reordenamos y llegamos a 

\[
\delta(\sqrt{-g}L)
=
\sqrt{-g}
\left[
\underbrace{\left(
\Rcal_{ab}
-\frac12g_{ab}L
-2\nabla^m\nabla^nP_{amnb}
\right)}_{E_{ab}}\delta g^{ab}
+
\nabla_j\delta v^j
\right]
\]

La variación de la acción queda

\begin{equation}
\boxed{
\delta A
=
\int_V\dd^Dx\,\sqrt{-g}\,E_{ab}\delta g^{ab}
+
\int_V\dd^Dx\,\sqrt{-g}\,\nabla_j\delta v^j
}
\label{eq:variacion-accion-final}
\end{equation}

donde

\[
E_{ab}
=
\Rcal_{ab}
-\frac12g_{ab}L
-2\nabla^m\nabla^nP_{amnb}
\]

\end{frame}



%================================================
\begin{frame}
	\frametitle{Identidad de Bianchi desde la covariancia}
	\small
	
	Para la teoría considerada por Padmanabhan,
	
	\[
	L=L(g^{ab},R_{abcd}),
	\]
	
	la variación de la acción puede escribirse como
	
	\[
	\delta A
	=
	\int_V d^Dx\,\sqrt{-g}
	\left(
	E_{ab}\,\delta g^{ab}
	+
	\nabla_a\Theta^a
	\right).
	\]
	
	Consideremos ahora una variación producida por un
	difeomorfismo infinitesimal generado por $\xi^a$:
	
	\[
	\delta_\xi g^{ab}
	=
	-\left(
	\nabla^a\xi^b+\nabla^b\xi^a
	\right)
	\]
	
	Como $E_{ab}=E_{ba}$, entonces 
	
	\[
	E_{ab}\delta_{\xi}g^{ab}=-2E_{ab}\nabla^b\xi^a.
	\]
	
	
	
\end{frame}

%================================================
\begin{frame}
	\frametitle{Identidad de Bianchi generalizada de Padmanabhan}
	\small
	
	Tomando $\xi^a$ con soporte compacto, los términos de
	frontera no contribuyen. Entonces,
	
	\begin{align*}
		0
		&=
		\int_V d^Dx\,\sqrt{-g}\,
		E_{ab}\,\delta_\xi g^{ab}
		=
		-2
		\int_V d^Dx\,\sqrt{-g}\,
		E_{ab}\nabla^a\xi^b
	\end{align*}
	
	Integramos por partes:
	
	\[
	-2E_{ab}\nabla^a\xi^b
	=
	-2\nabla^a(E_{ab}\xi^b)
	+
	2(\nabla^aE_{ab})\xi^b.
	\]
	
	El primer término es nuevamente una divergencia. Por tanto,
	
	\[
	0
	=
	2\int_V d^Dx\,\sqrt{-g}\,
	(\nabla^aE_{ab})\xi^b.
	\]
	
	Como $\xi^b(x)$ es arbitrario,
	
	\begin{equation}
		\boxed{
			\nabla^aE_{ab}\equiv0
		}
		\label{eq:Bianchi-Padmanabhan}
	\end{equation}
	
	Esta es la identidad de Bianchi generalizada asociada
	a la invariancia bajo difeomorfismos.
	
\end{frame}

%================================================
\begin{frame}
	\frametitle{Forma de la identidad en Lanczos--Lovelock}
	\small
	
	Para un lagrangiano general $L(g,R)$,
	
	\[
	E_{ab}
	=
	\mathcal{R}_{ab}
	-\frac12g_{ab}L
	-2\nabla^m\nabla^nP_{amnb},
	\]
	
	donde $\mathcal{R}_{ab} \equiv P_a{}^{cde}R_{bcde}$. En Lanczos--Lovelock se cumple además $\nabla_aP^{abcd}=0$, por lo que
	
	\[
	E_{ab}
	=
	\mathcal{R}_{ab}
	-\frac12g_{ab}L.
	\]
	
	Entonces, usando \Eq{eq:Bianchi-Padmanabhan},
	
	\[
	0
	=
	\nabla^aE_{ab}
	=
	\nabla^a\mathcal{R}_{ab}
	-\frac12\nabla_bL,
	\]
	
	y obtenemos
	
	\begin{equation}
		\boxed{
			\nabla^a\mathcal{R}_{ab}
			=
			\frac12\nabla_bL
		}
		\label{eq:Bianchi-LL-Rcal}
	\end{equation}
	
\end{frame}




%================================================
\begin{frame}
\frametitle{Caso Einstein--Hilbert: Recuperación del tensor de Einstein}

Recordemos que previamente habíamos encontrado que $\mathcal R_{ab}=R_{ab}$ para el caso de HE.

Como la conexión de Levi--Civita es compatible con la métrica, $\nabla_m g^{ab}=0$, se tiene
\[
\nabla_m P^{abcd}_{EH}=0.
\]

Por tanto,
\[
\nabla^m\nabla^nP_{amnb}=0.
\]

Así,
\[
\begin{aligned}
E_{ab}&=\mathcal R_{ab}-\frac12 g_{ab}L-2\nabla^m\nabla^nP_{amnb}
=R_{ab}-\frac12 g_{ab}R
\end{aligned}
\]

Por lo tanto,
\[
\boxed{
E_{ab}=G_{ab}
}
\qquad
\longrightarrow
\qquad
\nabla^a G_{ab} = 0
\]

Es decir, para $L=R$  y $I_{HE}$ en la generalización de Lanczos-Lovelock recuperamos el tensor de Einstein $G_{ab}$ junto a la identidad de Bianchi.
\end{frame}

%================================================
\begin{frame}
\frametitle{Término de frontera para Einstein--Hilbert}
\footnotesize

La expresión general $\delta v^j = 2P^{ibjd}\nabla_b\delta g_{di} - 2\delta g_{di}\nabla_cP^{ijcd}$ se reduce a 

\[
\delta v_{\mathrm{EH}}^{j}
=
\left(
g^{ij}g^{bd}-g^{id}g^{bj}
\right)
\nabla_b\delta g_{di}.
\]

Para escribirla en términos de \(\delta g^{ab}\), usamos $\delta g_{di} = -g_{da}g_{ic}\,\delta g^{ac}$.

La contracción del primer término es

\begin{align*}
g^{ij}g^{bd}\nabla_b\delta g_{di}
&=
-g^{ij}g^{bd}
\nabla_b
\left(
g_{da}g_{ic}\delta g^{ac}
\right)
=
-g^{ij}g^{bd}g_{da}g_{ic}
\nabla_b\delta g^{ac}
\\
&=
-\underbrace{g^{bd}g_{da}}_{\delta^b_a}
 \underbrace{g^{ij}g_{ic}}_{\delta^j_c}
\nabla_b\delta g^{ac}
=
-\delta^b_a\delta^j_c
\nabla_b\delta g^{ac}
=
-\nabla_a\delta g^{aj}.
\end{align*}

Análogamente, para el segundo término,

\begin{align*}
-g^{id}g^{bj}\nabla_b\delta g_{di}
&=
g^{id}g^{bj}g_{da}g_{ic}
\nabla_b\delta g^{ac}
=
\underbrace{g^{id}g_{da}}_{\delta^i_a}
g_{ic}\,
\underbrace{g^{bj}\nabla_b}_{\nabla^j}
\delta g^{ac}
=
g_{ac}\nabla^j\delta g^{ac}.
\end{align*}

En consecuencia,

\[
\boxed{
\delta v_{\mathrm{EH}}^{j}
=
g_{ab}\nabla^j\delta g^{ab}
-
\nabla_a\delta g^{aj}
}
\]

\end{frame}
%================================================\



%================================================
\begin{frame}
\frametitle{Constante cosmológica: ecuaciones y frontera}
\small

La expresión general del tensor de ecuaciones de campo es

\[
E_{ab}
=
\mathcal{R}_{ab}
-\frac12 g_{ab}L
-2\nabla^m\nabla^nP_{amnb}.
\]

Para \(L_{\Lambda}=-2\Lambda\), usando

\[
\mathcal{R}_{ab}^{(\Lambda)}=0,
\qquad
P_{amnb}^{(\Lambda)}=0,
\]

se obtiene

\begin{align*}
E_{ab}^{(\Lambda)}
&=
-\frac12g_{ab}(-2\Lambda)
=
\Lambda g_{ab}
\qquad
\longrightarrow
\qquad
\nabla^a E^{(\Lambda)}_{ab} = 0
\end{align*}

Por otro lado, gracias a que \(P_{\Lambda}^{abcd}=0\), el término de frontera es,

\[
\delta v_{\Lambda}^{j}
=
2P_{\Lambda}^{ibjd}\nabla_b\delta g_{di}
-
2\delta g_{di}\nabla_cP_{\Lambda}^{ijcd}=0
\]

Así, la constante cosmológica contribuye al término de volumen,
pero no genera ningún término de frontera.

\end{frame}



%================================================
\begin{frame}
\frametitle{Resumen de la sección}

\RoadmapCue{
Hemos aislado la contribución dinámica y el término de borde
que aparecen en $\delta I$.
}

\scriptsize

La variación de la acción toma la forma

\[
\delta A
=
\int_V \dd^D x\,\sqrt{-g}\,
E_{ab}\,\delta g^{ab}
+
\int_V \dd^D x\,\sqrt{-g}\,
\nabla_j\delta v^j
\]

con

\[
E_{ab}
=
\mathcal R_{ab}
-\frac12 g_{ab}L
-2\nabla^m\nabla^nP_{amnb}.
\]

El término de borde queda determinado por

\[
\delta v^j
=
2P^{ibjd}\nabla_b\delta g_{di}
-
2\delta g_{di}\nabla_cP^{ijcd}
\]


El único término que puede introducir derivadas de la métrica de orden superior a dos es $ -2\nabla^m\nabla^nP_{amnb}$. Para Einstein--Hilbert,

\[
P^{abcd}_{EH}
=
\frac12
\left(
g^{ac}g^{bd}-g^{ad}g^{bc}
\right),
\qquad
\nabla_mP^{abcd}_{EH}=0,
\]

y por tanto

\[
E_{ab}
=
R_{ab}-\frac12g_{ab}R
=
G_{ab} ;\qquad
\delta v_{\mathrm{EH}}^{j}
=
g_{ab}\nabla^j\delta g^{ab}
-
\nabla_a\delta g^{aj}
\]

\end{frame}



%================================================
\section{Condición de segundo orden}

\SetRoadmapStage{3}
%================================================
\begin{frame}
\frametitle{Por qué aparece el problema de cuarto orden}
\small

Recordemos que

\[
P^{abcd}
=
\pdv{L}{R_{abcd}}.
\]

Si $L$ depende de la curvatura, entonces $P^{abcd}$ también depende de la curvatura.

Pero

\[
R_{abcd}\sim \partial^2g+(\partial g)^2.
\]

Entonces,

\[
\nabla\nabla P
\sim
\nabla\nabla R
\sim
\partial^4g+\cdots.
\]

Por tanto, el término

\[
-2\nabla^m\nabla^nP_{amnb}
\]

puede introducir derivadas de cuarto orden de la métrica.

\end{frame}

%================================================
\begin{frame}
\frametitle{Condición clave}
\small

Para evitar derivadas mayores que segundo orden, imponemos

\[
\boxed{
\nabla_aP^{abcd}=0.
}
\]

Por las simetrías de $P^{abcd}$, esto implica divergencia nula en todos sus índices relevantes.

Entonces,

\[
\nabla^m\nabla^nP_{amnb}=0.
\]

Por tanto, el tensor de campo se reduce a

\[
\boxed{
E_{ab}
=
\Rcal_{ab}
-\frac12g_{ab}L.
}
\]

Con materia,

\[
\boxed{
\Rcal_{ab}
-\frac12g_{ab}L
=
\frac12T_{ab}.
}
\]

\end{frame}

%================================================
\begin{frame}
\frametitle{Resultado conceptual}
\small

La pregunta inicial era:

\begin{quote}
¿Qué lagrangianos construidos con métrica y curvatura generan ecuaciones de movimiento de segundo orden?
\end{quote}

La deducción muestra que el problema se reduce a buscar escalares $L$ tales que

\[
\nabla_aP^{abcd}=0,
\qquad
P^{abcd}
=
\pdv{L}{R_{abcd}}.
\]

\begin{alertblock}{Resultado}
Los lagrangianos que cumplen esta condición son precisamente los de Lanczos-Lovelock.
\end{alertblock}

La forma general es $L=\sum_m c_mL_m$, con
\[L_m\propto \delta^{a_1b_1\cdots a_mb_m}_{c_1d_1\cdots c_md_m}
R^{c_1d_1}{}_{a_1b_1} \cdots R^{c_md_m}{}_{a_mb_m}
\]

\end{frame}

%================================================
\section{Generalización con campo escalar}
\SetRoadmapStage{4}
%================================================

\begin{frame}

\frametitle{Generalización del lagrangiano}
\scriptsize

Hasta ahora hemos trabajado con un lagrangiano de tipo Lanczos--Lovelock puro, $L=L(g^{ab},R_{abcd})$. Ahora generalizamos a

\[
\boxed{
L=L(g^{ab},R_{abcd},\phi,\nabla_a\phi)
}
\]

Trabajaremos directamente con \(g^{ab}\), pues resulta más
conveniente cuando aparezcan objetos como

\[
\nabla^a\phi=g^{ab}\nabla_b\phi.
\]

Definimos

\begin{align}
P^{abcd}
&\equiv
\left(
\frac{\partial L}{\partial R_{abcd}}
\right)_{g^{mn},\phi,\nabla\phi}, \qquad
P_{ab} 
\equiv
\left(
\frac{\partial L}{\partial g^{ab}}
\right)_{R_{ijkl},\phi,\nabla\phi},
\\[1mm]
J^a
&\equiv
\left(
\frac{\partial L}{\partial(\nabla_a\phi)}
\right)_{g^{mn},R_{ijkl},\phi}, \qquad
F_\phi
\equiv
\left(
\frac{\partial L}{\partial\phi}
\right)_{g^{mn},R_{ijkl},\nabla\phi}.
\end{align}

\begin{block}{Cambio respecto al caso LL puro}
Como ahora derivamos respecto de la métrica contravariante
\(g^{ab}\), el objeto correspondiente lleva índices covariantes:
\(P_{ab}\).
\end{block}

\end{frame}


%================================================
\begin{frame}
\frametitle{Variación del lagrangiano generalizado}
\small

Con las definiciones anteriores,

\begin{equation}
\boxed{
\delta L
=
P_{ab}\delta g^{ab}
+
P^{abcd}\delta R_{abcd}
+
F_\phi\,\delta\phi
+
J^a\delta(\nabla_a\phi)
}
\label{eq:variacion-L-generalizado}
\end{equation}

La acción es

\[
S = \int_V \dd^Dx\,\sqrt{-g}\,L
\]

Por tanto, $\delta(\sqrt{-g}L) = \sqrt{-g}\,\delta L + L\,\delta\sqrt{-g}$ con $\delta\sqrt{-g} = -\frac12\sqrt{-g}\,g_{ab}\delta g^{ab}$, obtenemos
\[
\begin{aligned}
\delta(\sqrt{-g}L)
=
\sqrt{-g}\Big[
&P_{ab}\delta g^{ab}
+P^{abcd}\delta R_{abcd}
+F_\phi\delta\phi
+J^a\delta(\nabla_a\phi) -\frac12g_{ab}L\,\delta g^{ab}
\Big]
\end{aligned}
\]

Como \(\phi\) es un escalar, $\nabla_a\phi=\partial_a\phi$, y por ello $\boxed{\delta(\nabla_a\phi)=\nabla_a(\delta\phi)}$

Entonces,

\begin{equation}
\delta(\sqrt{-g}L)
=
\sqrt{-g}
\left[
P^{abcd}\delta R_{abcd}
+
\left(
P_{ab}-\frac12g_{ab}L
\right)\delta g^{ab}
+
F_\phi\delta\phi
+
J^a\nabla_a\delta\phi
\right]
\label{eq:variacion-densidad-generalizada}
\end{equation}

\end{frame}


%================================================
\begin{frame}
\frametitle{Separación de las variaciones}
\scriptsize

Para el sector gravitacional reutilizamos el resultado ya obtenido:

\begin{equation}
P^{abcd}\delta R_{abcd}
=
-\left(
\Rcal_{ab}
+
2\nabla^m\nabla^nP_{amnb}
\right)\delta g^{ab}
+
\nabla_a\delta v^a.
\label{eq:sector-gravitacional-generalizado}
\end{equation}

donde \(\delta v^a\) es el término de frontera definido previamente
en \Eq{eq:def-delta-v}. Ahora estudiamos $J^a\nabla_a(\delta\phi)$. Aplicando la regla del producto,

\[
\nabla_a(J^a\delta\phi)
=
(\nabla_aJ^a)\delta\phi
+
J^a\nabla_a(\delta\phi).
\]

Por tanto,

\begin{equation}
\boxed{
J^a\nabla_a(\delta\phi)
=
\nabla_a(J^a\delta\phi)
-
(\nabla_aJ^a)\delta\phi.
}
\label{eq:integracion-partes-escalar}
\end{equation}

Así, dentro de la acción,

\[
\int_V\dd^Dx\,\sqrt{-g}\,
J^a\nabla_a(\delta\phi)
=
\int_V\dd^Dx\,\sqrt{-g}
\left[
-(\nabla_aJ^a)\delta\phi
+
\nabla_a(J^a\delta\phi)
\right]
\]

\begin{block}{Estructura}
El primer término contribuye a la ecuación de campo de \(\phi\),
mientras que $\nabla_a(J^a\delta\phi)$ es una nueva contribución de frontera.
\end{block}

\end{frame}


%================================================
\begin{frame}
	\frametitle{Identidad generalizada para \(\Rcal_{ab}\)}
	\small
	
	Para una teoría general
	
	\[
	L=L(g^{ab},R_{abcd},\phi,u_a),
	\qquad
	u_a\equiv\nabla_a\phi,
	\]
	
	la variación del lagrangiano tiene la forma
	
	\[
	\delta L
	=
	P_{ab}\,\delta g^{ab}
	+
	P^{abcd}\,\delta R_{abcd}
	+
	J^a\,\delta u_a
	+
	F_\phi\,\delta\phi.
	\]
	
	Bajo un difeomorfismo generado por \(\xi^a\), $\delta_\xi g^{ab} = -2\nabla^{(a}\xi^{b)}$, mientras que, como \(u_a\) es un covector,	
	\[
	\delta_\xi u_a
	=
	\xi^c\nabla_cu_a
	+
	u_c\nabla_a\xi^c.
	\]
	
	Usando además las simetrías de \(P^{abcd}\), los términos
	proporcionales a \(\nabla_a\xi_b\) conducen a	
	\[
	-2P^{ab}
	+
	4\Rcal^{ab}
	+
	J^a u^b
	=
	0,
	\]	
	donde $\Rcal_{ab} \equiv P_a{}^{cde}R_{bcde}$. Por tanto,
	
	\begin{equation}
		\boxed{
			\Rcal_{ab} = \frac12 P_{ab} - \frac14 J_a u_b			
		}
		\label{eq:identidad-Rcal-general}
	\end{equation}
	
\end{frame}


%================================================
\begin{frame}
\frametitle{Variación completa de la acción}
\scriptsize

Reemplazando \Eq{eq:sector-gravitacional-generalizado} y
\Eq{eq:integracion-partes-escalar} en
\Eq{eq:variacion-densidad-generalizada},

\[
\begin{aligned}
\delta S
=
\int_V\dd^Dx\,\sqrt{-g}\Big[
&
-\left(
\Rcal_{ab}
+
2\nabla^m\nabla^nP_{amnb}
\right)\delta g^{ab}
+
\nabla_a\delta v^a
\\
&+
\left(
P_{ab}-\frac12g_{ab}L
\right)\delta g^{ab}
+
F_\phi\delta\phi
-
(\nabla_aJ^a)\delta\phi
+
\nabla_a(J^a\delta\phi)
\Big].
\end{aligned}
\]

con $\delta v^a = 2P^{ibjd}\nabla_b\delta g_{di} - 2\delta g_{di}\nabla_cP^{ijcd}$. Agrupando las variaciones independientes, definimos

\begin{align}
\boxed{
E_{ab}
=
P_{ab}
-\frac12g_{ab}L
-\Rcal_{ab}
-2\nabla^m\nabla^nP_{amnb}
}
\label{eq:Eab-generalizado}
\\[1mm]
\boxed{
E_\phi
=
F_\phi-\nabla_aJ^a
}
\label{eq:Ephi-generalizado}
\\[1mm]
\boxed{
\Theta^a
=
\delta v^a+J^a\delta\phi
}
\label{eq:theta-generalizado}
\end{align}

De modo que

\begin{equation}
\boxed{
\delta S
=
\int_V\dd^Dx\,\sqrt{-g}
\left[
E_{ab}\delta g^{ab}
+
E_\phi\delta\phi
+
\nabla_a\Theta^a
\right].
}
\label{eq:variacion-accion-generalizada-final}
\end{equation}

\end{frame}


%================================================
\begin{frame}
	\frametitle{Ecuación métrica como extensión de $L(g,R)$}
	\small
	
	Para $L=L(g^{ab},R_{abcd},\phi,u_a), u_a=\nabla_a\phi$, la identidad algebraica asociada al carácter escalar de $L$ se generaliza a $P_{ab}=2\Rcal_{ab}+ \frac12 J_{a}u_{b}$. Por tanto,	
	\begin{align*}
		E_{ab}
		&=
		P_{ab}
		-\frac12g_{ab}L
		-\Rcal_{ab}
		-2\nabla^m\nabla^nP_{amnb}
		\\
		&=
		\boxed{
			\Rcal_{ab}
			-\frac12g_{ab}L
			-2\nabla^m\nabla^nP_{amnb}
			+
			\frac12J_{a}u_{b}
		}
	\end{align*}
	
	Así, la estructura queda separada como
	
	\[
	\boxed{
		E_{ab}^{(g,R,\phi,\nabla\phi)}
		=
		E_{ab}^{(g,R)}
		+
		\frac12J_{a}u_{b}
	},
	\]
	
	entendiendo que $L$, $P^{abcd}$ y $\Rcal_{ab}$ son ahora
	los correspondientes a la teoría completa. Al apagar el sector escalar, $\phi\to0, u_a=\nabla_a\phi\to0$, se recupera inmediatamente
	
	\[
		E_{ab}^{(g,R)}
		=
		\Rcal_{ab}
		-\frac12g_{ab}L
		-2\nabla^m\nabla^nP_{amnb}
	\]
	
\end{frame}


%================================================
\begin{frame}
	\frametitle{Identidad de Bianchi en la teoría generalizada}
	\small
	
	Consideremos ahora nuestra extensión
	
	\[
	L=L(g^{ab},R_{abcd},\phi,\nabla_a\phi).
	\]
	
	La variación completa tiene la forma
	
	\[
	\delta A
	=
	\int_V d^Dx\,\sqrt{-g}
	\left[
	E_{ab}\delta g^{ab}
	+
	E_\phi\delta\phi
	+
	\nabla_a\Theta^a
	\right].
	\]
	
	Bajo un difeomorfismo generado por $\xi^a$, $\delta_\xi g^{ab} = -2\nabla^{a}\xi^{b}$, mientras que, por ser $\phi$ un escalar,
	
	\[
	\delta_\xi\phi
	=
	\xi^a\nabla_a\phi.
	\]
	
	Definiendo $u_a\equiv\nabla_a\phi$, la parte de volumen de la variación queda
	
	\[
	0
	=
	\int_V d^Dx\,\sqrt{-g}
	\left[
	-2E_{ab}\nabla^a\xi^b
	+
	E_\phi\,u_b\xi^b
	\right]
	\]
	\end{frame}

%================================================
\begin{frame}
	\frametitle{Identidad de Bianchi en la teoría generalizada}
	\small
	
	Integramos por partes el término métrico:
	
	\[
	-2E_{ab}\nabla^a\xi^b
	=
	-2\nabla^a(E_{ab}\xi^b)
	+
	2(\nabla^aE_{ab})\xi^b.
	\]
	
	Descartando la divergencia total por contribuir a la frontera,
	
	\[
	0
	=
	\int_V d^Dx\,\sqrt{-g}
	\left[
	2\nabla^aE_{ab}
	+
	E_\phi u_b
	\right]\xi^b.
	\]
	
	Como $\xi^b$ es arbitrario,
	
	\begin{equation}
		\boxed{
			2\nabla^aE_{ab}
			+
			E_\phi\nabla_b\phi
			\equiv0
		}
		\label{eq:Bianchi-generalizada-escalar}
	\end{equation}
	
	Esta identidad es \emph{off-shell}: no hemos impuesto todavía
	
	\[
	E_{ab}=0,
	\qquad
	E_\phi=0.
	\]
	
	Si la ecuación escalar se satisface,
	
	\[
	E_\phi=0
	\quad\Longrightarrow\quad
	\nabla^aE_{ab}=0.
	\]
	
	Para una teoría sin campo escalar se recupera inmediatamente
	la identidad de Padmanabhan.
	
\end{frame}

%================================================
\begin{frame}
	\frametitle{Forma útil para teorías con simetría de desplazamiento}
	\small
	
	En los términos que hemos estudiado no existe dependencia
	explícita en $\phi$. Por tanto,
	
	\[
	F_\phi\equiv\frac{\partial L}{\partial\phi}=0.
	\]
	
	Si definimos el momento escalar
	
	\[
	J^a\equiv
	\frac{\partial L}{\partial(\nabla_a\phi)},
	\]
	
	la ecuación escalar es
	
	\[
	E_\phi
	=
	-\nabla_aJ^a.
	\]
	
	Así, \Eq{eq:Bianchi-generalizada-escalar} puede escribirse como
	
	\begin{equation}
		\boxed{
			2\nabla^aE_{ab}
			-
			(\nabla_aJ^a)u_b
			\equiv0.
		}
		\label{eq:Bianchi-J}
	\end{equation}
	
	Esta forma será especialmente útil para los sectores
	$\mathcal{L}_0$ y $Q_\beta$.
	
\end{frame}

%================================================
\begin{frame}
	\frametitle{Identidad de Bianchi en términos de los momentos}
	\small
	
	Partimos de
	
	\[
	2\nabla^aE_{ab}
	+
	E_\phi\nabla_b\phi
	\equiv0,
	\]
	
	con
	
	\[
	E_{ab}
	=
	P_{ab}
	-\frac12g_{ab}L
	-\Rcal_{ab}
	-2\nabla^m\nabla^nP_{amnb},
	\qquad
	E_\phi=F_\phi-\nabla_aJ^a.
	\]
	
	Sustituyendo y usando $\nabla^ag_{ab}=0$,
	
	\begin{equation}
	\boxed{
		2\nabla^aP_{ab}
		-\nabla_bL
		-2\nabla^a\Rcal_{ab}
		-4\nabla^a\nabla^m\nabla^nP_{amnb}
		+
		\left(F_\phi-\nabla_aJ^a\right)u_b
		\equiv0
		\label{Eq:BianchiGeneralizado}
	}
	\end{equation}
	
	Para simetría de desplazamiento, $F_\phi=0$:
	
	\[
	\boxed{
		2\nabla^aP_{ab}
		-\nabla_bL
		-2\nabla^a\Rcal_{ab}
		-4\nabla^a\nabla^m\nabla^nP_{amnb}
		-
		(\nabla_aJ^a)u_b
		\equiv0
	}
	\]
	
\end{frame}


%================================================
\begin{frame}
	\frametitle{Ecuaciones de campo}
	\scriptsize
	
	Ignorando o cancelando adecuadamente los términos de frontera,
	la condición \(\delta S=0\) para variaciones arbitrarias implica
	
	\[
	\boxed{
		E_{ab}=0,
		\qquad
		E_\phi=0.
	}
	\]
	
	Por tanto, la ecuación métrica es
	
	\begin{equation}
		\boxed{
			P_{ab}
			-\frac12g_{ab}L
			-\Rcal_{ab}
			-2\nabla^m\nabla^nP_{amnb}
			=
			0.
		}
		\label{eq:campo-metrico-generalizado}
	\end{equation}
	
	Mientras que para el campo escalar,
	\begin{align}
		F_\phi-\nabla_aJ^a&=0 \qquad \longrightarrow \qquad
		\boxed{
			\frac{\partial L}{\partial\phi}
			-
			\nabla_a
			\left(
			\frac{\partial L}
			{\partial(\nabla_a\phi)}
			\right)
			=0}
		\label{eq:euler-lagrange-escalar}
	\end{align}
	
	Esta es precisamente la ecuación de Euler--Lagrange covariante
	para el campo escalar.
	
	\begin{block}{Caso sin dependencia explícita en \(\phi\)}
		Si $L=L(g^{ab},R_{abcd},\nabla_a\phi)$, entonces \(F_\phi=0\) y la ecuación escalar se reduce a
		\[
		\boxed{\nabla_aJ^a=0.}
		\]
		Esto refleja la simetría de desplazamiento
		\(\phi\rightarrow\phi+\text{cte}\).
	\end{block}
	
\end{frame}


%================================================
\section{Estudio de casos}
\SetRoadmapStage{5}
\subsection{Gauss - Bonet}
\begin{frame}{Gauss - Bonet}
    
\end{frame}

\subsection{Lagrangiano especial}
%================================================
\begin{frame}
\frametitle{Acción concreta: separación de los sectores}
\small

Consideramos la acción

\[
\boxed{
I_{\mathrm{EQT}}
=
\frac{1}{16\pi G}
\int \dd^3x\,\sqrt{-g}\,
\left(\mathcal{L}_0+Q_\beta\right)
}
\]

con

\[
\mathcal{L}_0 = R+\frac{2}{L^2} -\alpha_1(\nabla\phi)^2,
\qquad
Q_\beta = L^2\beta_0
\left[
3R_{ab}\nabla^a\phi\nabla^b\phi - (\nabla\phi)^2R
\right]
\]

Por ahora trabajaremos únicamente con \(\mathcal{L}_0\). Los sectores $R$ y $\frac{2}{L^2}$ ya fueron desarrollados anteriormente, por lo que resta estudiar

\[
\mathcal{L}_{\mathrm{cin}}
= -\alpha_1(\nabla\phi)^2
\]

Definimos $u_a\equiv\nabla_a\phi$, $u^a=g^{ab}u_b$, de manera que

\[
\boxed{
\mathcal{L}_{\mathrm{cin}}
=
-\alpha_1g^{ab}u_a u_b
}
\]

\end{frame}


%================================================
\begin{frame}
\frametitle{Momentos de \(\mathcal{L}_{\mathrm{cin}}\)}
\small

Como no existe dependencia en el tensor de Riemann,

\[
\boxed{
P_{\mathrm{cin}}^{abcd} = \frac{\partial \mathcal{L}_{\mathrm{cin}}}
{\partial R_{abcd}}
= 0
}
\]

y, por tanto,

\[
\Rcal^{\mathrm{cin}}_{ab}=0,
\qquad
\nabla^m\nabla^n
P^{\mathrm{cin}}_{amnb}=0.
\]

Para el momento métrico,

\[
P^{\mathrm{cin}}_{ab}
=
\frac{\partial \mathcal{L}_{\mathrm{cin}}}{\partial g^{ab}},
\]

tenemos

\[
\delta_g \mathcal{L}_{\mathrm{cin}}
=
-\alpha_1u_a u_b\,\delta g^{ab},
\]

de modo que

\[
\boxed{
P^{\mathrm{cin}}_{ab} = -\alpha_1u_a u_b.
}
\]

\end{frame}


%================================================
\begin{frame}
\frametitle{Momentos escalares de \(\mathcal{L}_{\mathrm{cin}}\)}
\small

Para el momento asociado a \(\nabla_a\phi\),

\[
J^a_{\mathrm{cin}}
=
\frac{\partial \mathcal{L}_{\mathrm{cin}}}{\partial u_a}
\]

Variando respecto de \(u_a\) y teniendo en cuenta la simetría de $g^{ab}$,

\begin{align*}
\delta_u \mathcal{L}_{\mathrm{cin}}
&=
-\alpha_1 g^{ab}
\left(
\delta u_a\,u_b
+
u_a\,\delta u_b
\right)=
-2\alpha_1u^a\,\delta u_a
\end{align*}

Por tanto,

\[
J^a_{\mathrm{cin}}
=
-2\alpha_1u^a
=
-2\alpha_1\nabla^a\phi
\]

Además, \(\mathcal{L}_{\mathrm{cin}}\) no depende explícitamente de \(\phi\), por lo que

\[
F_\phi^{\mathrm{cin}}
=
\frac{\partial \mathcal{L}_{\mathrm{cin}}}{\partial\phi}
= 0
\]

Así, los momentos del sector cinético son

\[
\boxed{
P_{\mathrm{cin}}^{abcd}=0,
\qquad
P^{\mathrm{cin}}_{ab}=-\alpha_1u_a u_b,
\qquad
J^a_{\mathrm{cin}}=-2\alpha_1u^a,
\qquad
F_\phi^{\mathrm{cin}}=0
}
\]

\end{frame}


%================================================
\begin{frame}
\frametitle{Variación directa de \(\mathcal{L}_{\mathrm{cin}}\)}
\small

Por tanto,

\[
\boxed{
\delta \mathcal{L}_{\mathrm{cin}}
=
P^{\mathrm{cin}}_{ab}\delta g^{ab}
+
J^a_{\mathrm{cin}}\delta u_a.
}
\]

Como $u_a=\nabla_a\phi$ y \(\phi\) es un escalar, $\delta u_a = \delta(\nabla_a\phi) = \nabla_a(\delta\phi)$. Luego,

\[
\delta \mathcal{L}_{\mathrm{cin}}
=
-\alpha_1u_a u_b\,\delta g^{ab}
-
2\alpha_1u^a\nabla_a(\delta\phi)
\]


Ahora,

\[
\delta(\sqrt{-g}\,\mathcal{L}_{\mathrm{cin}})
=
\sqrt{-g}\,\delta \mathcal{L}_{\mathrm{cin}}
+
\mathcal{L}_{\mathrm{cin}}\delta\sqrt{-g}.
\]

Usando $\delta\sqrt{-g} = -\frac12\sqrt{-g}\, g_{ab}\delta g^{ab}$, obtenemos

\begin{align*}
\delta(\sqrt{-g}\,\mathcal{L}_{\mathrm{cin}})
=
\sqrt{-g}\Bigg[
&
\left(
P^{\mathrm{cin}}_{ab}
-\frac12g_{ab}\mathcal{L}_{\mathrm{cin}}
\right)\delta g^{ab}
+
J^a_{\mathrm{cin}}\delta u_a
\Bigg]
\end{align*}

Sustituyendo los momentos,

\[
\delta(\sqrt{-g}\,\mathcal{L}_{\mathrm{cin}})
=
\sqrt{-g}
\left[
\left(
-\alpha_1u_a u_b
+
\frac{\alpha_1}{2}g_{ab}u^cu_c
\right)\delta g^{ab}
-
2\alpha_1u^a\nabla_a\delta\phi
\right]
\]

\end{frame}


%================================================
\begin{frame}
\frametitle{Integración por partes del término escalar}
\small

Para el último término usamos $\nabla_a(u^a\delta\phi) = (\nabla_a u^a)\delta\phi + u^a\nabla_a\delta\phi$. Por tanto,

\[
u^a\nabla_a\delta\phi
=
\nabla_a(u^a\delta\phi)
-
(\nabla_a u^a)\delta\phi.
\]

Entonces,

\begin{align*}
\delta(\sqrt{-g}\,\mathcal{L}_{\mathrm{cin}})
=
\sqrt{-g}\Bigg[
&
\left(
-\alpha_1u_a u_b
+
\frac{\alpha_1}{2}g_{ab}u^cu_c
\right)\delta g^{ab}
\\
&
+
2\alpha_1(\nabla_a u^a)\delta\phi
-
2\alpha_1\nabla_a(u^a\delta\phi)
\Bigg]
\end{align*}

Identificamos así

\[
\boxed{
E^{\mathrm{cin}}_{ab} = -\alpha_1u_a u_b + \frac{\alpha_1}{2}g_{ab}u^cu_c,
\qquad
E^{\mathrm{cin}}_{\phi} = 2\alpha_1\nabla_a u^a = 2\alpha_1\nabla_a\nabla^a\phi,
}
\]

y el término de frontera

\[
\boxed{
\delta v^a_{\mathrm{cin}}
=
-2\alpha_1\nabla^a\phi\,\delta\phi.
}
\]

\end{frame}


%================================================
\begin{frame}
\frametitle{Reconstrucción del sector \(\mathcal{L}_0\)}
\small

Recordemos

\[
\mathcal{L}_0
=
R+\frac{2}{L^2}-\alpha_1(\nabla\phi)^2.
\]

Como los sectores EH y cosmológico ya son conocidos,

\[
E^{\mathrm{EH}}_{ab}=G_{ab},
\qquad
E^\Lambda_{ab}
=
-\frac{1}{L^2}g_{ab},
\]

Por tanto, definimos

\[
E^{(0)}_{ab}
=
E^{\mathrm{EH}}_{ab}
+
E^\Lambda_{ab}
+
E^{\mathrm{cin}}_{ab},
\]

\[
E^{(0)}_\phi
=
E^{\mathrm{cin}}_\phi.
\]

Para la frontera,

\[
\delta v^{a(0)}
=
\delta v^a_{\mathrm{EH}}
+
\delta v^a_{\mathrm{cin}},
\]

pues $\delta v^a_\Lambda=0$.

\end{frame}


%================================================
\begin{frame}
\frametitle{Resultado final para \(\mathcal{L}_0\)}
\small

Reemplazando las contribuciones,

\begin{equation}
\boxed{
E^{(0)}_{ab}
=
G_{ab}
-\frac{1}{L^2}g_{ab}
-\alpha_1\nabla_a\phi\nabla_b\phi
+
\frac{\alpha_1}{2}
g_{ab}(\nabla\phi)^2
}
\label{eq:Eab-h0}
\end{equation}

y

\begin{equation}
\boxed{
E^{(0)}_\phi
=
2\alpha_1\nabla_a\nabla^a\phi
=
2\alpha_1\Box\phi
}
\label{eq:Ephi-h0}
\end{equation}

El término de frontera total es

\begin{equation}
\boxed{
\delta v^{a(0)}
=
g_{bc}\nabla^a\delta g^{bc}
-
\nabla_b\delta g^{ab}
-
2\alpha_1\nabla^a\phi\,\delta\phi
}
\label{eq:frontera-h0}
\end{equation}

Así,

\[
\boxed{
\delta I_{\mathcal{L}_0}
=
\frac{1}{16\pi G}
\int\dd^3x\,\sqrt{-g}
\left[
E^{(0)}_{ab}\delta g^{ab}
+
E^{(0)}_\phi\delta\phi
+
\nabla_a\delta v^{a(0)}
\right]
}
\]

\end{frame}


%================================================
\begin{frame}
\frametitle{Sector de acoplamiento \(Q_\beta\)}
\small

Consideramos ahora

\[
\boxed{
Q_\beta
=
L^2\beta_0
\left[
3R_{ab}\nabla^a\phi\nabla^b\phi
-
(\nabla\phi)^2R
\right].
}
\]

Para simplificar la notación definimos

\[
u_a\equiv\nabla_a\phi,
\qquad
u^a=g^{ab}u_b,
\qquad
X\equiv u^a u_a=(\nabla\phi)^2.
\]

Además,

\[
R_{bd}=g^{ac}R_{abcd},
\qquad
R=g^{ac}g^{bd}R_{abcd}.
\]

Por tanto,

\begin{align*}
Q_\beta
&=
L^2\beta_0
\left[
3g^{ac}g^{bm}g^{dn}
-
g^{ac}g^{bd}g^{mn}
\right]
R_{abcd}u_m u_n
\\[1mm]
&=
L^2\beta_0
\left[
3g^{ac}u^bu^d
-
Xg^{ac}g^{bd}
\right]R_{abcd}.
\end{align*}

Calcularemos los cuatro momentos

\[
P_\beta^{abcd},
\qquad
P^{(\beta)}_{ab},
\qquad
J_\beta^a,
\qquad
F_\phi^{(\beta)}.
\]

\end{frame}


%================================================
\begin{frame}
\frametitle{Momento de curvatura \(P_\beta^{abcd}\)}
\small

De la forma anterior,

\[
P_\beta^{abcd}\delta R_{abcd}
=
L^2\beta_0
\left[
3g^{ac}u^bu^d
-
Xg^{ac}g^{bd}
\right]
\delta R_{abcd}.
\]

Como \(\delta R_{abcd}\) es anstisimétrico en los pares $(a, b)$ y $(c, d)$,
debemos proyectar cada coeficiente sobre esas mismas simetrías.

Para el primer término, aplicando 

\[
g^{ac}u^bu^d\,\delta R_{abcd}
=
\frac14
\left(
g^{ac}u^bu^d
-g^{bc}u^au^d
-g^{ad}u^bu^c
+g^{bd}u^au^c
\right)
\delta R_{abcd}.
\]

Mientras que

\[
g^{ac}g^{bd}\delta R_{abcd}
=
\frac12
\left(
g^{ac}g^{bd}
-
g^{ad}g^{bc}
\right)\delta R_{abcd}.
\]

Por tanto,

\begin{equation}
\boxed{
\begin{aligned}
P_\beta^{abcd}
=
\frac{L^2\beta_0}{4}
\Big[
&3\big(
g^{ac}u^bu^d
-g^{bc}u^au^d
-g^{ad}u^bu^c
+g^{bd}u^au^c
\big)
\\
&-2(\nabla\phi)^2
\big(
g^{ac}g^{bd}
-
g^{ad}g^{bc}
\big)
\Big].
\end{aligned}
}
\label{eq:Pbeta-curvatura}
\end{equation}

\end{frame}


%================================================
\begin{frame}
\frametitle{Momento métrico \(P^{(\beta)}_{ab}\)}
\small

Partimos de la expresión completamente desarrollada

\[
Q_\beta
=
L^2\beta_0
\left[
3g^{ac}g^{bm}g^{dn}
-
g^{ac}g^{bd}g^{mn}
\right]
R_{abcd}\,u_m u_n,
\]

donde $u_a\equiv\nabla_a\phi, X\equiv u^a u_a=(\nabla\phi)^2$. Para calcular \(P^{(\beta)}_{ab}\), variamos únicamente las
métricas, manteniendo \(R_{abcd}\) y \(u_a\) fijos:

\[
\delta_g Q_\beta
=
P^{(\beta)}_{pq}\,\delta g^{pq}.
\]

Así,

\[
\begin{aligned}
\delta_g Q_\beta
=
L^2\beta_0
\Big\{
3\Big[
&\delta g^{ac}g^{bm}g^{dn}
+g^{ac}\delta g^{bm}g^{dn}
+g^{ac}g^{bm}\delta g^{dn}
\Big]
\\
-\Big[
&\delta g^{ac}g^{bd}g^{mn}
+g^{ac}\delta g^{bd}g^{mn}
+g^{ac}g^{bd}\delta g^{mn}
\Big]
\Big\}
R_{abcd}u_m u_n .
\end{aligned}
\]

Como $\delta g^{pq}=\delta g^{qp}$, todo coeficiente \(A_{pq}\) puede reemplazarse por su parte simétrica:

\[
\boxed{
A_{pq}\delta g^{pq} = A_{(pq)}\delta g^{pq}
}
\]

\end{frame}

%================================================
\begin{frame}
\frametitle{Variación métrica: primer grupo}
\scriptsize

Consideremos primero los tres términos multiplicados por \(3\).

Para el primero,

\begin{align*}
\delta g^{ac}g^{bm}g^{dn}
R_{abcd}u_m u_n
&=
u^bu^dR_{abcd}\,\delta g^{ac} = u^iu^jR_{piqj}\,\delta g^{pq}
\end{align*}

Para el segundo usamos \(g^{ac}R_{abcd}=R_{bd}\):

\begin{align*}
g^{ac}\delta g^{bm}g^{dn}
R_{abcd}u_m u_n
&=
R_{bd}\,u^d u_m\,\delta g^{bm}=
\frac12
\left(
u_q u^iR_{pi}
+
u_p u^iR_{qi}
\right)
\delta g^{pq}
\end{align*}

Análogamente,

\begin{align*}
g^{ac}g^{bm}\delta g^{dn}
R_{abcd}u_m u_n
&=
R_{bd}\,u^b u_n\,\delta g^{dn}=
\frac12
\left(
u_q u^iR_{pi}
+
u_p u^iR_{qi}
\right)
\delta g^{pq}.
\end{align*}

Por tanto, los tres términos juntos dan

\begin{equation}
\boxed{
\begin{aligned}
&
\Big(
\delta g^{ac}g^{bm}g^{dn}
+
g^{ac}\delta g^{bm}g^{dn}
+
g^{ac}g^{bm}\delta g^{dn}
\Big)
R_{abcd}u_m u_n
\\[1mm]
&\qquad=
\left[
u^iu^jR_{piqj}
+
u_q u^iR_{pi}
+
u_p u^iR_{qi}
\right]\delta g^{pq}.
\end{aligned}
}
\label{eq:Pbeta-grupo-positivo}
\end{equation}

\end{frame}

%================================================
\begin{frame}
\frametitle{Variación métrica: segundo grupo}
\scriptsize

Ahora estudiamos

\[
\Big(
\delta g^{ac}g^{bd}g^{mn}
+
g^{ac}\delta g^{bd}g^{mn}
+
g^{ac}g^{bd}\delta g^{mn}
\Big)
R_{abcd}u_m u_n
\]

Para el primer término,

\begin{align*}
\delta g^{ac}g^{bd}g^{mn}
R_{abcd}u_m u_n
&=
X\,g^{bd}R_{abcd}\,\delta g^{ac}=
X R_{pq}\delta g^{pq}
\end{align*}

Para el segundo,

\begin{align*}
g^{ac}\delta g^{bd}g^{mn}
R_{abcd}u_m u_n
&=
X\,g^{ac}R_{abcd}\,\delta g^{bd}=
X R_{pq}\delta g^{pq}
\end{align*}

Finalmente,

\begin{align*}
g^{ac}g^{bd}\delta g^{mn}
R_{abcd}u_m u_n
&=
R\,u_m u_n\,\delta g^{mn}
=
R\,u_pu_q\,\delta g^{pq}
\end{align*}

Por tanto,

\begin{equation}
\boxed{
\begin{aligned}
&
\Big(
\delta g^{ac}g^{bd}g^{mn}
+
g^{ac}\delta g^{bd}g^{mn}
+
g^{ac}g^{bd}\delta g^{mn}
\Big)
R_{abcd}u_m u_n
\\
&\qquad=
\left[
2X R_{pq}
+
R\,u_pu_q
\right]\delta g^{pq}
\end{aligned}
}
\label{eq:Pbeta-grupo-negativo}
\end{equation}

\end{frame}

%================================================
\begin{frame}
\frametitle{Momento métrico \(P^{(\beta)}_{ab}\): resultado}
\small

Sustituyendo \Eq{eq:Pbeta-grupo-positivo} y
\Eq{eq:Pbeta-grupo-negativo},

\[
\begin{aligned}
\delta_g Q_\beta
=
L^2\beta_0
\Big\{
&3\left[
u^iu^jR_{piqj}
+
u_q u^iR_{pi}
+
u_p u^iR_{qi}
\right]
\\
&-
\left[
2XR_{pq}
+
Ru_pu_q
\right]
\Big\}
\delta g^{pq}
\end{aligned}
\]

Comparando con

\[
\delta_g Q_\beta
=
P^{(\beta)}_{pq}\delta g^{pq},
\]

obtenemos

\begin{equation}
\boxed{
\begin{aligned}
P^{(\beta)}_{pq}
=
L^2\beta_0
\Big[
&3u^iu^jR_{piqj}
+3u_q u^iR_{pi}
+3u_p u^iR_{qi}
\\
&-2XR_{pq}
-Ru_pu_q
\Big].
\end{aligned}
}
\label{eq:Pbeta-metrico-expandido}
\end{equation}

Finalmente, renombrando \(p,q\rightarrow a,b\),

\begin{equation}
\boxed{
P^{(\beta)}_{ab}
=
L^2\beta_0
\left[
3R_{acbd}u^cu^d
+
6u^c u_{(a}R_{b)c}
-
2XR_{ab}
-
Ru_a u_b
\right].
}
\label{eq:Pbeta-metrico}
\end{equation}

\end{frame}


%================================================
\begin{frame}
\frametitle{Momento escalar \(J_\beta^a\)}
\small

Para

\[
J_\beta^a
=
\frac{\partial Q_\beta}{\partial u_a},
\qquad
u_a\equiv\nabla_a\phi,
\]

partimos de

\[
Q_\beta
=
L^2\beta_0
\left(
3R_{ab}u^a u^b
-
Ru^a u_a
\right)
\]

Al variar únicamente respecto de \(u_a\),

\begin{align*}
\delta_u Q_\beta
=
L^2\beta_0
\Big[
&3R_{bd}
\left(
\delta u^b\,u^d
+
u^b\delta u^d
\right)
-
R\left(
\delta u^m\,u_m
+
u^m\delta u_m
\right)
\Big]
\end{align*}

Como la métrica se mantiene fija en esta variación,

\[
u^a=g^{ab}u_b
\qquad\Longrightarrow\qquad
\boxed{
\delta u^a=g^{ab}\delta u_b
}
\]

Por tanto,

\begin{align*}
\delta_u Q_\beta
=
L^2\beta_0
\Big[
&3R_{bd}
\left(
g^{bm}\delta u_m\,u^d
+
u^b g^{dm}\delta u_m
\right)
-
R\left(
g^{mn}\delta u_n\,u_m
+
u^m\delta u_m
\right)
\Big]
\end{align*}

\end{frame}

%================================================
\begin{frame}
\frametitle{Momento escalar \(J_\beta^a\): resultado}
\scriptsize

Reescribimos todos los términos en función de la misma
variación \(\delta u_m\):

\begin{align*}
\delta_u Q_\beta
=
L^2\beta_0
\Big[
&3R_{bd}g^{bm}u^d
+
3R_{bd}u^b g^{dm}
-
R\,g^{mn}u_n
-
Ru^m
\Big]\delta u_m
\end{align*}

Ahora, $g^{mn}u_n=u^m$, y, usando la simetría del tensor de Ricci, $R_{bd}=R_{db}$, tenemos

\[
R_{bd}g^{bm}u^d
=
R^m{}_d u^d,
\]

mientras que

\[
R_{bd}u^b g^{dm}
=
R_b{}^m u^b
=
R^m{}_b u^b.
\]

Por tanto, ambos términos son iguales y

\begin{align*}
\delta_u Q_\beta
&=
L^2\beta_0
\left[
6R^m{}_b u^b
-
2Ru^m
\right]\delta u_m
\end{align*}

Comparando con $\delta_u Q_\beta = J_\beta^m\delta u_m$, obtenemos

\begin{equation}
\boxed{
J_\beta^a
=
L^2\beta_0
\left(
6R^a{}_b\nabla^b\phi
-
2R\nabla^a\phi
\right)
}
\label{eq:Jbeta}
\end{equation}

Finalmente, como \(Q_\beta\) no depende explícitamente de \(\phi\),

\begin{equation}
\boxed{
F_\phi^{(\beta)} = 0
}
\label{eq:Fbeta}
\end{equation}

\end{frame}


%================================================
\begin{frame}
	\frametitle{Contracción \(\Rcal^{(\beta)}_{(ab)}\)}
	\scriptsize
	
	Para \(Q_\beta\), ya obtuvimos
	
	\[
	P^{(\beta)}_{ab}
	=
	L^2\beta_0
	\left[
	3R_{acbd}u^cu^d
	+
	6u^cu_{(a}R_{b)c}
	-
	2XR_{ab}
	-
	Ru_au_b
	\right],
	\]
	
	y
	
	\[
	J^{(\beta)}_a
	=
	L^2\beta_0
	\left(
	6R_{ac}u^c
	-
	2Ru_a
	\right).
	\]
	
	Usamos ahora directamente la identidad \Eq{eq:identidad-Rcal-general}: $\Rcal^{(\beta)}_{(ab)} = \frac12P^{(\beta)}_{ab} - \frac14J^{(\beta)}_{(a}u_{b)}$.
	
	El primer término es
	
	\begin{align*}
		\frac12P^{(\beta)}_{ab}
		=
		L^2\beta_0
		\Big[
		&
		\frac32R_{acbd}u^cu^d
		+
		3u_{(a}R_{b)c}u^c
		-XR_{ab}
		-\frac12Ru_au_b
		\Big]
	\end{align*}
	
	Mientras que
	
	\[
	\frac14J^{(\beta)}_{(a}u_{b)}
	=
	L^2\beta_0
	\left[
	\frac32u_{(a}R_{b)c}u^c
	-
	\frac12Ru_au_b
	\right]
	\]
	
	Restando, los términos \(Ru_au_b\) se cancelan y obtenemos
	
	\begin{equation}
		\boxed{
			\Rcal^{(\beta)}_{(ab)}
			=
			L^2\beta_0
			\left[
			\frac32R_{acbd}u^cu^d
			+
			\frac32u_{(a}R_{b)c}u^c
			-
			XR_{ab}
			\right].
		}
		\label{eq:Rcalbeta-sim}
	\end{equation}
	
	Es exactamente el resultado obtenido mediante
	la contracción directa
	\(P^{(\beta)}_a{}^{cde}R_{bcde}\), solo que este método es mucho más tedioso de hacer.
	
\end{frame}


%================================================
\begin{frame}
\frametitle{Término \(-2\nabla\nabla P_\beta\)}
\scriptsize

A partir de \Eq{eq:Pbeta-curvatura}, la contribución simétrica es

\[
-2\nabla_m\nabla_n
P^{(\beta)}_{(a}{}^{mn}{}_{b)}.
\]

Después de bajar los índices correspondientes y usar
\(\nabla_a g_{bc}=0\), obtenemos

\begin{equation}
\boxed{
\begin{aligned}
-2\nabla_m\nabla_n
P^{(\beta)}_{(a}{}^{mn}{}_{b)}
=
L^2\beta_0\Bigg\{
-\frac34\Big[
&
\nabla_m\nabla_a(u^m u_b)
+
\nabla_m\nabla_b(u^m u_a)
\\
&+
\nabla_b\nabla_n(u_a u^n)
+
\nabla_a\nabla_n(u_b u^n)
\Big]
\\
&+
\frac32\Box(u_a u_b)
+
\frac32g_{ab}\nabla_m\nabla_n(u^m u^n)
\\
&+
\nabla_a\nabla_bX
-
g_{ab}\Box X
\Bigg\}.
\end{aligned}
}
\label{eq:nablanabla-Pbeta}
\end{equation}

donde

\[
X=(\nabla\phi)^2,
\qquad
\Box\equiv\nabla^c\nabla_c.
\]

\end{frame}

%================================================
\begin{frame}
\frametitle{Variación de la acción \(I_{Q_\beta}\)}
\small

El sector correspondiente a \(Q_\beta\) es

\[
I_{Q_\beta}
=
\frac{1}{16\pi G}
\int\dd^3x\,\sqrt{-g}\,Q_\beta.
\]

Su variación es

\[
\delta I_{Q_\beta}
=
\frac{1}{16\pi G}
\int\dd^3x\,
\delta(\sqrt{-g}\,Q_\beta).
\]

Usando $\delta\sqrt{-g} = -\frac12\sqrt{-g}\,g_{ab}\delta g^{ab}$, tenemos

\[
\begin{aligned}
\delta I_{Q_\beta}
=
\frac{1}{16\pi G}
\int\dd^3x\,\sqrt{-g}
\Big[
&
P_\beta^{abcd}\delta R_{abcd}
\\
&+
\left(
P^{(\beta)}_{ab}
-\frac12g_{ab}Q_\beta
\right)\delta g^{ab}
+
J_\beta^a\nabla_a\delta\phi
\Big].
\end{aligned}
\]

Para el sector escalar,

\[
J_\beta^a\nabla_a\delta\phi
=
-(\nabla_aJ_\beta^a)\delta\phi
+
\nabla_a(J_\beta^a\delta\phi).
\]

\end{frame}

%================================================
\begin{frame}
\frametitle{Ecuaciones asociadas a \(Q_\beta\)}
\small

Usando el resultado general del sector de curvatura,

\[
P_\beta^{abcd}\delta R_{abcd}
=
-
\left[
\Rcal^{(\beta)}_{(ab)}
+
2\nabla_m\nabla_n
P^{(\beta)}_{(a}{}^{mn}{}_{b)}
\right]\delta g^{ab}
+
\nabla_a\delta v^a_{P,\beta},
\]

la variación adopta la forma

\[
\boxed{
\delta I_{Q_\beta}
=
\frac{1}{16\pi G}
\int\dd^3x\,\sqrt{-g}
\left[
E^{(\beta)}_{ab}\delta g^{ab}
+
E^{(\beta)}_\phi\delta\phi
+
\nabla_a\Theta^a_{(\beta)}
\right].
}
\]

donde

\begin{equation}
\boxed{
E^{(\beta)}_{ab}
=
P^{(\beta)}_{ab}
-\frac12g_{ab}Q_\beta
-\Rcal^{(\beta)}_{(ab)}
-2\nabla_m\nabla_n
P^{(\beta)}_{(a}{}^{mn}{}_{b)}.
}
\label{eq:Eab-beta}
\end{equation}

y

\begin{equation}
\boxed{
E^{(\beta)}_\phi
=
-\nabla_aJ_\beta^a
=
-2L^2\beta_0
\nabla_a
\left(
3R^a{}_b\nabla^b\phi
-
R\nabla^a\phi
\right).
}
\label{eq:Ephi-beta}
\end{equation}

\end{frame}

%================================================
\begin{frame}
\frametitle{Término de frontera de \(Q_\beta\)}
\small

La contribución de frontera queda

\[
\Theta^j_{(\beta)}
=
\delta v^j_{P,\beta}
+
J_\beta^j\delta\phi,
\]

con

\[
\delta v^j_{P,\beta}
=
2P_\beta^{ibjd}\nabla_b\delta g_{di}
-
2\delta g_{di}\nabla_cP_\beta^{ijcd}.
\]

Por tanto,

\begin{equation}
\boxed{
\begin{aligned}
\Theta^j_{(\beta)}
={}&
2P_\beta^{ibjd}\nabla_b\delta g_{di}
-
2\delta g_{di}\nabla_cP_\beta^{ijcd}
\\
&+
2L^2\beta_0
\left(
3R^j{}_b\nabla^b\phi
-
R\nabla^j\phi
\right)\delta\phi.
\end{aligned}
}
\label{eq:frontera-Qbeta}
\end{equation}

\begin{block}{Resultado}
El sector \(Q_\beta\) produce contribuciones tanto a la
ecuación métrica como a la ecuación escalar y al término
de frontera.
\end{block}

\end{frame}



%================================================
\begin{frame}
	\frametitle{Chequeo de Bianchi para $\mathcal{L}_0+Q_\beta$}
	\small
	
	Para
	
	\[
	L
	=
	R+\frac{2}{L^2}
	-\alpha_1 X
	+
	L^2\beta_0
	\left(
	3R_{ab}u^au^b-RX
	\right),
	\qquad
	u_a=\nabla_a\phi,
	\]
	
	la identidad generalizada exige $2\nabla^aE_{ab}+E_\phi u_b=0$.	Usando los resultados obtenidos para cada sector,
	
	\begin{align*}
		2\nabla^aE_{ab}
		={}&
		\underbrace{2\nabla^aG_{ab}}_{0}
		- \underbrace{\frac{2}{L^2}\nabla^ag_{ab}}_{0}
		- 2\alpha_1(\Box\phi)\,u_b
		+ L^2\beta_0 \nabla_a \left(6R^a{}_c u^c-2Ru^a\right)u_b
	\end{align*}
	
	mientras que la ecuación escalar es
	
	\[
	E_\phi
	=
	2\alpha_1\Box\phi
	-
	L^2\beta_0
	\nabla_a
	\left(
	6R^a{}_c u^c-2Ru^a
	\right).
	\]
	
	Por tanto,
	
	\begin{align*}
		2\nabla^aE_{ab}+E_\phi u_b
		={}&
		-2\alpha_1(\Box\phi)u_b
		+
		L^2\beta_0
		\nabla_a
		\left(
		6R^a{}_c u^c-2Ru^a
		\right)u_b
		\\
		&+
		2\alpha_1(\Box\phi)u_b
		-
		L^2\beta_0
		\nabla_a
		\left(
		6R^a{}_c u^c-2Ru^a
		\right)u_b
		=0
	\end{align*}
	
\end{frame}





%================================================
\begin{frame}
\frametitle{Ecuaciones de campo de la teoría completa}
\small

La densidad lagrangiana total es

\[
\mathcal{L}
=
\mathcal{L}_0+Q_\beta.
\]

Por linealidad de la variación,

\[
E_{ab}
=
E^{(0)}_{ab}
+
E^{(\beta)}_{ab},
\]

\[
E_\phi
=
E^{(0)}_\phi
+
E^{(\beta)}_\phi,
\]

y

\[
\Theta^a
=
\Theta^a_{(0)}
+
\Theta^a_{(\beta)}.
\]

Por tanto, a partir de \Eq{eq:Eab-h0}, \Eq{eq:Ephi-h0}, \Eq{eq:frontera-h0}, \Eq{eq:Eab-beta}, \Eq{eq:Ephi-beta} y \Eq{eq:frontera-Qbeta}, las ecuaciones de campo son

\[
\boxed{
E^{(0)}_{ab}+E^{(\beta)}_{ab}=0
}
\qquad
\boxed{
E^{(0)}_\phi+E^{(\beta)}_\phi=0
}
\]

\end{frame}



%================================================
\section{Discusión de resultados}
\SetRoadmapStage{5}
%================================================
%================================================
\begin{frame}
	\frametitle{Condición de Lanczos--Lovelock: $\nabla_a P^{abcd}$}
	\small
	
	En Lanczos--Lovelock se cumple la condición estructural $\nabla_a P^{abcd}=0$. Para nuestro modelo,
	
	\[
	\mathcal L
	=
	R+\frac{2}{L^2}
	-\alpha_1(\nabla\phi)^2
	+
	Q_\beta,
	\]
	
	ni la constante cosmológica ni el término cinético contribuyen a \(P^{abcd}\), pues no dependen explícitamente de \(R_{abcd}\): $P_{\Lambda}^{abcd}=0, P_{\mathrm{cin}}^{abcd}=0$. Además, $\nabla_a P_{\mathrm{EH}}^{abcd}=0$ por compatibilidad métrica, \(\nabla_a g^{bc}=0\). Por tanto,
	
	\[
	\boxed{
		\nabla_aP^{abcd}
		=
		\nabla_aP_\beta^{abcd}
	}
	\]
	
	Partimos entonces del tensor ya obtenido:
	
	\begin{equation*}
		\begin{aligned}
			P_\beta^{abcd}
			=
			\frac{L^2\beta_0}{4}
			\Big[
			&3\big(
			g^{ac}u^bu^d
			-g^{bc}u^au^d
			-g^{ad}u^bu^c
			+g^{bd}u^au^c
			\big)
			\\
			&-2(\nabla\phi)^2
			\big(
			g^{ac}g^{bd}
			-
			g^{ad}g^{bc}
			\big)
			\Big]
		\end{aligned}
	\end{equation*}
	con
	
	\[
	u_a=\nabla_a\phi,
	\qquad
	X=u^a u_a.
	\]
	
\end{frame}


%================================================
\begin{frame}
	\frametitle{Divergencia de $P_\beta^{abcd}$}
	\small
	
	Como \(\nabla_a g^{bc}=0\), la derivada covariante actúa únicamente
	sobre \(u^a\) y \(X\):
	
	\begin{align*}
		\nabla_aP_\beta^{abcd}
		=
		\frac{L^2\beta_0}{4}
		\Bigg\{
		3\Big[
		&
		u^d\nabla^c u^b
		+u^b\nabla^c u^d
		-g^{bc}u^d\nabla_a u^a
		-g^{bc}u^a\nabla_a u^d
		\\
		&
		-u^c\nabla^d u^b
		-u^b\nabla^d u^c
		+g^{bd}u^c\nabla_a u^a
		+g^{bd}u^a\nabla_a u^c
		\Big]
		\\
		&
		-2
		\left(
		g^{bd}\nabla^c X
		-
		g^{bc}\nabla^d X
		\right)
		\Bigg\}.
	\end{align*}
	
	Como \(u_a=\nabla_a\phi\),
	
	\[
	\nabla_a u_b=\nabla_b u_a,
	\qquad
	\nabla_a u^a=\Box\phi,
	\]
	
	y
	
	\[
	\nabla^dX
	=
	\nabla^d(u^a u_a)
	=
	2u^a\nabla_a u^d.
	\]
	
	Además, 
	$u^b\nabla^c u^d-u^b\nabla^d u^c=0$.
	
\end{frame}


%================================================
\begin{frame}
	\frametitle{Expresión general para $\nabla_aP_\beta^{abcd}$}
	\small
	
	Usando las identidades anteriores, la divergencia se reduce a
	
	\[
	\boxed{
		\begin{aligned}
			\nabla_aP_\beta^{abcd}
			=
			\frac{L^2\beta_0}{4}
			\Bigg\{
			3\Big[
			&
			u^d
			\left(
			\nabla^c u^b-g^{bc}\Box\phi
			\right)
			\\
			&
			-u^c
			\left(
			\nabla^d u^b-g^{bd}\Box\phi
			\right)
			\Big]
			\\
			&
			+\frac12
			\left(
			g^{bc}\nabla^dX
			-
			g^{bd}\nabla^cX
			\right)
			\Bigg\}.
		\end{aligned}
	}
	\]
	
	Por tanto, para nuestro lagrangiano completo,
	
	\[
	\boxed{
		\nabla_aP^{abcd}
		=
		\nabla_aP_\beta^{abcd}
	}.
	\]
	
	A diferencia del caso Lanczos--Lovelock, la expresión anterior
	\emph{no se anula identicamente}.
	
	Para verificarlo sobre la solución concreta, basta encontrar una
	sola componente no nula.
	
\end{frame}


%================================================
\begin{frame}
	\frametitle{Evaluación sobre el ansatz de la solución}
	\small
	
	Consideramos
	
	\[
	ds^2
	=
	-f(r)dt^2
	+
	\frac{dr^2}{f(r)}
	+
	r^2d\varphi^2,
	\qquad
	d\phi=p\,d\varphi.
	\]
	
	Entonces,
	
	\[
	u_a=(0,0,p),
	\qquad
	u^a=
	\left(
	0,0,\frac{p}{r^2}
	\right),
	\qquad
	\boxed{
		X=\frac{p^2}{r^2}
	}.
	\]
	
	Además, $\Box\phi=0$ por la relación en el ansatz. Pero, $\nabla_a u_b\neq0$, pues, por ejemplo,
	
	\[
	\nabla_r u_\varphi
	=
	-\Gamma^\varphi_{r\varphi}u_\varphi
	=
	-\frac{p}{r}
	\]
	
	Para demostrar que \(\nabla_aP^{abcd}\neq0\), tomamos la componente	
	\[
	(b,c,d)=(t,t,r)
	\]
	
	Como $u^t=u^r=0, g^{tr}=0$, la expresión general se reduce a	
	\[
	\boxed{
		\nabla_aP_\beta^{attr}
		=
		\frac{L^2\beta_0}{8}
		g^{tt}\nabla^rX
	}
	\]
	
\end{frame}


%================================================
\begin{frame}
	\frametitle{No se cumple la condición de Lanczos--Lovelock}
	\scriptsize
	
	Como
	
	\[
	g^{tt}=-\frac{1}{f(r)},
	\qquad
	\text{y}
	\qquad
	\nabla^rX
	=
	g^{rr}\partial_rX
	=
	f(r)\,
	\partial_r\left(\frac{p^2}{r^2}\right)
	=
	-\frac{2f(r)p^2}{r^3},
	\]
	
	obtenemos
	
	\begin{align*}
		\nabla_aP_\beta^{attr}
		&=
		\frac{L^2\beta_0}{8}
		\left(-\frac1f\right)
		\left(-\frac{2fp^2}{r^3}\right)
		=
		\boxed{
			\frac{L^2\beta_0p^2}{4r^3}
		}
	\end{align*}
	
	Por tanto,
	
	\[
	\boxed{
		\nabla_aP^{attr}\neq0
		\quad\Longrightarrow\quad
		\nabla_aP^{abcd}\neq0.
	}
	\]
	
	Notemos que \(f(r)\) se cancela completamente. En consecuencia,
	sustituir específicamente la solución
	
	\[
	f(r)
	=
	\frac{
		r^2/L^2-\lambda
	}{
		1+\beta_0p^2L^2/r^2
	}
	\]
	
	no puede producir la cancelación. Para \(r\) finito, $\nabla_aP^{attr}=0$ solo si
	
	\[
	\boxed{
		\beta_0=0
		\qquad\text{o}\qquad
		p=0,
		\qquad\text{o}\qquad
		r\longrightarrow \infty
	}
	\]
	
	Los dos primeros corresponden a los casos triviales en los que desaparece
	la interacción \(Q_\beta\).
	
\end{frame}
%================================================






\begin{frame}{Referencias}
\footnotesize
\begin{thebibliography}{9}

\bibitem{Padmanabhan2010}
T. Padmanabhan,
\textit{Thermodynamical Aspects of Gravity: New insights},
Reports on Progress in Physics \textbf{73}, 046901 (2010),
arXiv:0911.5004.

\bibitem{PadmanabhanKothawala2013}
T. Padmanabhan and D. Kothawala,
\textit{Lanczos--Lovelock models of gravity},
Physics Reports \textbf{531}, 115--171 (2013),
arXiv:1302.2151.

\bibitem{Wald1993}
R. M. Wald,
\textit{Black hole entropy is the Noether charge},
Physical Review D \textbf{48}, R3427--R3431 (1993).

\end{thebibliography}
\end{frame}


\end{document}